\section{Answer Reduction}
\label{sec:ans}

In this section we show how to transform a normal form verifier $\verifier =
(\sampler, \decider)$ into an ``answer reduced'' normal form verifier
$\verifier^\ar = (\sampler^\ar, \decider^\ar)$ such that the values of the
associated nonlocal games are directly related, yet the answer-reduced
verifier's decision runtime is only polylogarithmic in the answer length of the
original verifier (the answer-reduced verifier's sampling runtime remains
polynomially related to the sampling runtime of $\sampler$).
The polylogarithmic dependence is achieved by composing a probabilistically
checkable proof (PCP) with the oracularized verifier given in \Cref{sec:oracle}.
This step generalizes the answer reduction technique of~\cite[Part V]{NW19}.

Given index~$n \in \N$, the answer reduced verifier $\verifier^\ar$ simulates
the oracularization~$\verifier^\ora$ of $\verifier$ on index $n$.
To do so, it first samples questions~$x$ and~$y$ using the oracularized sampler
$\sampler^\ora$ and distributes them to the players, who compute answers~$a$
and~$b$.
Let us suppose that the first player is assigned the $\oracle$ role, and parse
their question and answer as pairs $x = (x_\alice, x_\bob)$ and $a = (a_\alice,
a_\bob)$, while the second player is an ``isolated'' player receiving the
question $x_\alice$ and responding with answer $b$.
Instead of executing the decider $\decider^\ora$ on the answers $(a, b)$, the
verifier $\verifier^\ar$ asks the first player to compute a PCP $\Pi$ of
$\decider(n, x_\alice, x_\bob, a_\alice, a_\bob) = 1$, and the second player to
compute an encoding $g_b$ of answer $b$.
$\verifier^\ar$ then requests randomly chosen locations of the proof $\Pi$ and
the encoding $g_b$, and executes the PCP verifier on the players' answers.
By the soundness of the PCP, $\verifier^\ar$ accepts with high probability only
if the player's answers satisfy $\decider(n, x_\alice, x_\bob, a_\alice, a_\bob)
= 1$ and $b = a_\alice$.

There are several challenges that arise when implementing answer reduction.
One challenge, already encountered in~\cite{NW19}, is that we need to ensure the
PCP $\Pi$ computed by the first player can be cross-tested against the encoding
$g_b$ computed by the second player (who doesn't know the entire structure of
the PCP $\Pi$).
This was handled in~\cite{NW19} by using a special type of PCP called a
\emph{probabilistically checkable proof of proximity (PCPP)}, which allows one
to efficiently check that a \emph{specific} string $x$ is a satisfying
assignment to a Boolean formula $\varphi$, as opposed to simply checking that
$\varphi$ is satisfiable.
In a PCPP, an encoding of the specific string $x$ is provided separately from
the proof of satisfiability.
The answer reduction scheme of~\cite{NW19} was able to use an ``off-the-shelf''
PCPP in a relatively black-box fashion to handle this.

In our answer reduction scheme, however, there is a further requirement: we need
the question distribution of the answer reduced verifier to be conditionally
linear.
This is necessary to maintain the invariant that the verifier after each step of
the compression procedure (introspection, answer reduction, parallel repetition)
is a normal form verifier.
Unfortunately, simulating the question distributions of off-the-shelf PCPPs with
conditionally linear distributions can be quite cumbersome.
Instead, we design a bespoke PCP verifier for the protocol whose question
distribution is more easily seen to be conditionally linear.

This section is organized as follows.
We start with some preliminaries on formulas and encodings in
Section~\ref{sec:ar-tms}.
In Section~\ref{sec:cook-levin} we show how to use the Cook-Levin reduction to
reduce the Bounded Halting problem for deciders to a succinct satisfiability
problem called Succinct-3SAT.
Following this, in Section~\ref{sec:succinct-deciders}, we reduce the
Succinct-3SAT instance to an instance of a related problem called Succinct
Decoupled 5SAT, which is easier to use in our answer reduction step.
Then in Section~\ref{sec:pcp-cktval-new} we introduce a PCP for Succinct
Decoupled 5SAT.
The verifier for the PCP expects a proof consisting of the evaluation tables of
low-degree polynomials, including the low-degree encodings of the players'
answers $a$ and $b$.
In Section~\ref{sec:ld-compiler} we provide the definition of a normal-form
verifier $\verifier^\ar$ that executes the composition of $\verifier^\ora$ with
the PCP verifier from Section~\ref{sec:pcp-cktval-new}.
In Section~\ref{sec:ar-completeness} we show completeness of the construction
and analyze its complexity.
In Section~\ref{sec:ar-soundness} we prove soundness.

\subsection{Circuit preliminaries}
\label{sec:ar-tms}

Recall the definitions pertaining to Turing machines fromSection~\ref{sec:tms}. In this section it is useful to model
computation using \emph{Boolean circuits}. A Boolean circuit is
a network of Boolean gates, connected by directed ``wires.'' Incoming
wires encode the input to the circuit, and outgoing wires encode the
output. For the most part, we will assume that the output of the
circuit is a single bit, corresponding to a single output wire. The
in-degree of each gate (the ``fan-in'')
is restricted be at most 2, but there is no restriction on
the out-degree (the ``fan-out''). The \emph{size} of a circuit is the total
number of gates and wires that it contains. For a more detailed description, see Section 4.3 of~\cite{Pap94}.

\begin{remark}[Plugging integers into circuits]\label{rem:plugging-in-integers}
  Let $\circuit$ be a circuit with a single input of length~$n$.
  Inputs to~$\circuit$ are strings $x \in \{0, 1\}^n$.
  In this section, we will also allow~$\circuit$ to receive inputs $a \in \{0,
  1, \ldots, 2^n-1\}$.
  In doing so, we use the convention that a number~$a$ between~$0$ and~$2^{n}-1$
  is interpreted as its $n$-digit binary encoding $\binary{n}(a)$
when provided as input to a set of
  $n$ single-bit wires.
  In other words, $\circuit(a) = \circuit(x)$, where $x = \binary{n}(a)$.

  More generally, if the circuit $\circuit$ has~$k$ different inputs of length
  $n_1, \ldots, n_k$, then we can evaluate it on inputs $a_1 \in \{0, 1, \ldots,
  2^{n_1}-1\}, \ldots, a_k \in \{0, 1, \ldots, 2^{n_k}-1\}$ as follows:
  \begin{equation*}
    \circuit(a_1, \ldots, a_k) = \circuit(x_1, \ldots, x_k)\;,
  \end{equation*}
  where $x_1 = \binary{n_1}(a_1), \ldots, x_k = \binary{n_k}(a_k)$.
\end{remark}

A 3SAT formula is a Boolean formula in conjunctive normal form in which at most
three literals appear in each clause.
More precisely, $\varphi$ is a 3SAT formula on $N$ variables $x_1, x_2, \ldots,
x_N$ if it has the form $\bigwedge_{j=1}^m C_j$ and each clause $C_j$ is the
disjunction of at most three literals, where a literal is either a variable
$x_i$ or its negation $\neg x_i$.
We use $x_i^o$ to denote the literal $x_i$ if $o = 1$ and $\neg x_i$ if $o = 0$.

\begin{definition}[Succinct description of 3SAT formulas]
  \label{def:succinct-formulas}
  Let $N = 2^n$, and let $\varphi$ be a 3SAT formula on~$N$ variables named
  $x_0, \ldots, x_{N-1}$.
  Let $\circuit$ be a Boolean circuit with 3 inputs of length~$n$ and three
  single-bit inputs.
  Then $\circuit$ is a \emph{succinct description of $\varphi$} if for each
  $i_1, i_2, i_3 \in \{0, 1, \ldots, N-1\}$ and $o_1, o_2, o_3 \in \{0, 1\}$,
  \begin{equation}\label{eq:circuit-val-is-one}
    \circuit(i_1, i_2, i_3, o_1, o_2, o_3) = 1
  \end{equation}
  if and only if $x_{i_1}^{o_1} \lor x_{i_2}^{o_2} \lor x_{i_3}^{o_3}$ is a
  clause in $\varphi$.
  In Equation~\eqref{eq:circuit-val-is-one}, we use the notation from
  Remark~\ref{rem:plugging-in-integers}.
\end{definition}
  
\begin{definition}[Succinct-3SAT problem]
  The \emph{Succinct-3SAT} problem is the language containing encodings of
  circuits~$\circuit$ in which $\circuit$ is a succinct description of a
  satisfiable 3SAT formula~$\varphi$.
\end{definition}



The Tseitin transformation is a mapping from circuits to Boolean
formulas. The following summarizes its properties; for an explicit
construction, see \cite[Section 3.8]{NW19}.
\begin{definition}
  \label{def:tseitin}
  Let $\circuit$ be a Boolean circuit with $n$ inputs and size $s$. Then the corresponding \emph{Tseitin formula} $\formula$ is a
  Boolean formula on $n+s$ variables with the property that for all $x
  \in \{0,1\}^n$, $\circuit(x) = 1$ if and only if there exists $w \in
\{0,1\}^s$ such that $\formula(x, s) = 1$. The formula $\F$ has size
$O(s)$ and a description of it is computable from $\circuit$ in time $O(s)$.
\end{definition}

\begin{definition}
  \label{def:formula-arithmetization}
  Let $\formula$ be a Boolean formula over $m$ variables. The arithmetization
  $\formula_{\mathrm{arith}}$ over $\F_q$ is a function
  $\formula_{\mathrm{arith}}:\F_q^{m'}\rightarrow \F_q$ such that
  \begin{equation}
    \label{eq:farith}
    \forall x \in \{0, 1\}^{m}\;,
    \quad \formula_{\mathrm{arith}}(x) = \formula(x)\;.
  \end{equation}
  
\end{definition}

\begin{proposition}
  \label{prop:tseitin-arith-degree}
  Let $\circuit$ be a circuit on $n$ inputs with size $s$. Then the
  arithmetization $\formula_{\mathrm{arith}}$ of its Tseitin formula
  is a polynomial over $\F_q$ on $n +s$ variables with individual
  degree at most $2$ in each variable. For admissible field sizes $q$,
  a description of the polynomial
  $\formula_{\mathrm{arith}}$ can be computed in $\poly(s, \log q)$ time
  it can be evaluated at a specific point $z \in \F_q^{n+s}$ in time
  $\poly(s, \log q)$.
\end{proposition}
\begin{proof}
    The properties follow by inspecting the construction in Definition
    3.27 and 3.28 of \cite{NW19}, together with
    \Cref{lem:efficient_arithmetic}. In particular, the degree bound
    follows by observing that every variable in
    $\formula$ occurs at most twice, and therefore the arithmetization
    $\formula_{\mathrm{arith}}$ has individual degree $2$.
\end{proof}

\subsection{A Cook-Levin theorem for bounded deciders}
\label{sec:cook-levin}





\begin{definition}[Bounded Halting problem]
	The \emph{$k$-input Bounded Halting} problem is the language
  $\BoundedHalting_k$ containing the set of tuples $(\alpha, T, z_1,\ldots,z_k)$
  where $\alpha$ is the description of a $k$-input Turing machine, $T \in \N$,
  $z_1, \ldots, z_k \in \{0,1\}^*$, and $\cal{M}_\alpha$ accepts input $(z_1,
  \ldots, z_k)$ in at most $T$ time steps.
\end{definition}

We begin by defining natural encodings of a decider's tape alphabet and set of
states.

\begin{definition}[Decider encodings]
  Let~$\decider$ be a decider with tape alphabet $\Gamma = \{0, 1, \sqcup\}$ and
  set of states~$K$.
  We will write $\mathrm{enc}_{\Gamma}: \Gamma \cup \{\square\} \rightarrow \{0,
  1\}^2$ for the function which encodes the elements of~$\Gamma$, and a special
  ``$\square$'' symbol described below, as length-two binary strings in the
  following manner:
  \begin{alignat*}{2}
    & \mathrm{enc}_{\Gamma}(0) = 00\;,\quad
    && \mathrm{enc}_{\Gamma}(1) = 01\;,\\
    & \mathrm{enc}_{\Gamma}(\sqcup) = 10\;,\quad
    && \mathrm{enc}_{\Gamma}(\square) = 11\;.
  \end{alignat*}
  In addition, we write $\mathrm{enc}_K : K \rightarrow \{0, 1\}^\kappa$ for
  some arbitrary fixed $\kappa$-bit encoding of the elements of~$K$, where
  $\kappa = \lceil \log(|K|) \rceil$.
\end{definition}

Now, we give the main result of this section.
It states that any decider~$\decider$ can be converted into a circuit~$\circuit$
which succinctly represents a 3SAT formula~$\varphi_{\mathrm{3SAT}}$ that
carries out the time~$T$ computation of~$\decider$.
In addition, $\circuit$ is extremely small---size $\poly \log(T)$ rather than
$\poly(T)$.

\begin{proposition}[Succinct representation of deciders]
  \label{prop:standard-succinct-sat}
  There is an algorithm with the following properties.
  Let $\decider$ be a decider, let $n$, $T$,~$\qlen$, and $\sigma$ be integers
  with $\qlen \leq T$ and $|\decider| \leq \sigma$, and let~$x$ and~$y$ be
  strings of length at most~$\qlen$.
  Then on input $(\decider, n, T, \qlen, \sigma, x, y)$, the algorithm outputs
  a circuit~$\circuit$ on $3m + 3$ inputs which succinctly describes a 3SAT
  formula $\varphi_{\mathrm{3SAT}}$ on $M = 2^m$ variables.
  Furthermore, $\varphi_{\mathrm{3SAT}}$ has the following property:
  \begin{itemize}
  \item For all $a, b \in \{0, 1\}^{2T}$, there exists a $c \in \{0, 1\}^{M-4T}$
    such that $w = (a, b, c)$ satisfies $\varphi_{\mathrm{3SAT}}$ if and only if
    there exist $a_{\mathrm{prefix}}, b_{\mathrm{prefix}} \in \{0, 1\}^*$ of
    lengths~$\ell_a, \ell_b \leq T$, respectively, such that
    \begin{equation*}
      a = \mathrm{enc}_\Gamma(a_{\mathrm{prefix}}, \sqcup^{T - \ell_a})
      \quad
      \text{and}
      \quad
      b = \mathrm{enc}_\Gamma(b_{\mathrm{prefix}}, \sqcup^{T - \ell_b})
    \end{equation*}
    and $\decider$ accepts $(n, x, y, a_{\mathrm{prefix}}, b_{\mathrm{prefix}})$
    in time~$T$.
  \end{itemize}
  Finally, the following statements hold:
  \begin{enumerate}
  \item The parameter $m$ controlling the number of inputs to the circuit
    depends only on $T$ and $\sigma$, and $m(T,\sigma) = O(\log(T) +
    \log(\sigma))$,
  \item $\circuit$ has at most $s(n,T,\qlen,\sigma)=\poly(\log(n), \log(T),
    \qlen, \sigma)$ gates,
  \item The algorithm runs in time $\poly(\log(n), \log(T), \qlen,\sigma)$,
  \item Furthermore, explicit values for $m(T, \sigma)$ and
    $s(n, T, \qlen, \sigma)$ can be computed in time polynomial in
    $n,\log(T),\qlen,\sigma$.
  \end{enumerate}
\end{proposition}

Proposition~\ref{prop:standard-succinct-sat} is essentially the standard fact
that Succinct-3SAT is an $\mathsf{NEXP}$-complete language, i.e.\ that every
nondeterministic computation which takes time~$2^n$ can be represented as a
Succinct-3SAT instance of size only~$\poly(n)$.
However, it has several peculiarities which requires us to prove it from scratch
rather than simply appealing to the $\mathsf{NEXP}$-completeness of
Succinct-3SAT.
First, we require that the coordinates of~$a$ and~$b$ embed into~$w$ not
randomly but as its lexicographically first coordinates (for reasons that are
explained below in Section~\ref{sec:succinct-deciders}).
Second, we need explicit bounds on how quantities such as the size of~$\circuit$
relate to quantities such as $\sigma$, an upper bound on the description length
of $\decider$ in bits.


To prove Proposition~\ref{prop:standard-succinct-sat}, we follow the standard
proof that Succinct-3SAT is $\mathsf{NEXP}$-complete as presented
in~\cite{Pap94}.
This proof observes that the Cook-Levin reduction, which is used to show that
3SAT is $\mathsf{NP}$-complete, produces a 3SAT instance whose clauses follow
such a simple pattern that they can be described succinctly using an
exponentially-smaller circuit.
One key difference in our proof is that we will apply the Cook-Levin reduction
directly to the $5$-input Turing machine~$\decider$, which by \
Section~\ref{sec:tms} has~$7$ tapes; traditional proofs such as the one
in~\cite{Pap94} would first convert~$\decider$ to a single-tape Turing
machine~$\decider_{\mathrm{single}}$, and then apply the Cook-Levin reduction
for single-tape Turing machines to~$\decider_{\mathrm{single}}$.
Though this adds notational overhead to our proof, it allows us to more easily
track of which variables in $\varphi_{\mathrm{3SAT}}$ correspond to the
strings~$a$ and~$b$ (see Proposition~\ref{prop:standard-succinct-sat} to see
what these refer to).

\begin{proof}[Proof of Proposition~\ref{prop:standard-succinct-sat}]
  The Cook-Levin reduction considers the \emph{execution tableau} of~$\decider$
  when run for time~$T$.
  The execution tableau contains, for each time $t \in \{1, \ldots, T\}$,
  variables describing the state of~$\decider$ and the contents of each of its
  tape cells at time~$t$.
  More formally, it consists of the following three sets of variables.
  \begin{enumerate}
  \item \label{item:cell-variables} For each time $t \in \{1, \ldots, T\}$, tape
    $i \in \{1, \ldots, 7\}$, and tape position $j \in \{0, 1,\ldots, T+1\}$,
    the tableau contains two Boolean-valued variables
    \begin{equation*}
      c_{t, i, j} = (c_{t, i, j, 1}, c_{t, i, j, 2}) \in \{0, 1\}^2
    \end{equation*}
    which are supposed to correspond to the contents of the $j$-th tape cell on
    tape~$i$ at time~$t$ according to $\mathrm{enc}_{\Gamma}(\cdot)$.
    The variables with $j \in \{0, T+1\}$ do not correspond to any cell on the
    tape; rather, the $j=0$ variables correspond to the left-boundary of the
    tape, and the $j=T+1$ variables correspond to the right-boundary of the
    first~$T$ cells on the tape.
    These are expected to always contain the special boundary symbol
    ``$\square$'', i.e.\ $c_{t, i, j}$ should be equal to
    $\mathrm{enc}_{\Gamma}(\square)$ whenever $j \in \{0, T+1\}$.
    As we will see below, it is convenient to define these so that for each~$t
    \in \{1, \ldots, T\}$, $i\in \{1, \ldots, 7\}$, and $j \in \{1, \ldots,
    T\}$, the variable $c_{t, i, j}$ also has a variable to its left $c_{t, i,
      j-1}$ and to its right $c_{t, i, j+1}$.
  \item \label{item:head-variables} For each time $t \in \{1, \ldots, T\}$, tape
    $i \in \{1, \ldots, 7\}$, and tape position $j \in \{0, 1,\ldots, T+1\}$,
    the tableau contains Boolean-valued variables $h_{t, i, j} \in \{0, 1\}$
    which are supposed to indicate whether the $i$-th tape head is in cell~$j$
    at time~$t$.
    For the boundary cells $j \in \{0, T+1\}$, we expect that $h_{t, i, j} = 0$
    for all $t \in \{1, \ldots, T\}$ and $i \in \{1, \ldots, 7\}$.
  \item For each time $t\in \{1, \ldots, T\}$, the tableau contains $\kappa$
    Boolean-valued variables
    \begin{equation*}
      s_t = (s_{t, 1}, \ldots, s_{t, \kappa}) \in \{0, 1\}^\kappa
    \end{equation*}
    which are supposed to correspond to the state of $\decider$ at time~$t$
    according to $\mathrm{enc}_K(\cdot)$.
  \end{enumerate}
  Finally, we let $\mathcal{V}$ denote the set of all of these variables.
  In other words,
  \begin{equation*}
    \mathcal{V} = \{c_{t, i, j, k}\}_{t, i, j, k} \cup \{h_{t, i, j}\}_{t, i, j}
    \cup \{s_{t, k}\}_{t, k}.
  \end{equation*}
  In total, the number of variables in the execution tableau is given by
  \begin{equation}\label{eq:tableau-variable-count}
    |\mathcal{V}|
    = O\big(T^2 + T \cdot \log(|K|)\big)
    = O\big(T^2 + T \cdot \log(|\decider|)\big)\;.
  \end{equation}
  The first term in Equation~\eqref{eq:tableau-variable-count}
  corresponds to the tape cell encodings~$c_{t, i, j}$ and $h_{t, i, j}$,
  and the second term corresponds to the Turing machine state encodings~$s_t$.
  The second equality uses the fact that $|K| \leq |\decider|$.

  As stated above, we expect the variables in~$\mathcal{V}$ to correspond to
  some time-$T$ execution of the decider~$\decider$.
  However, in general these are just arbitrary $\{0, 1\}$-valued variables.
  We now describe a set of constraints placed on these variables which, if
  satisfied, ensure they \emph{do} indeed correspond to some time-$T$ execution
  of~$\decider$.
  These constraints will be split into two categories: (i) the constraints
  corresponding to the boundary, which ensure that the $t=1$ variables are
  initialized to a valid starting configuration and the $j \in \{0, T+1\}$
  variables are set according to Items~\ref{item:cell-variables}
  and~\ref{item:head-variables}, and (ii) the constraints corresponding to the
  execution of~$\decider$, which ensure that the variables at each time~$(t+1)$
  follow from the variables at time~$t$ according to the computation
  of~$\decider$.
  We start with the boundary constraints, which are simple enough to be
  described with a 3SAT formula.

  \begin{definition}
    The \emph{boundary formula} $\varphi_{\mathrm{Boundary}}$ is the 3SAT
    formula on the variables $\mathcal{V}$ described as follows.
    Let $o = \mathrm{enc}_K({\mathsf{start}}) \in \{0, 1\}^\kappa$, where
    ${\mathsf{start}} \in K$ is the start state of~$\decider$.
    For the $t= 1$ boundary, $\varphi_{\mathrm{Boundary}}$ contains the
    following set of clauses.
    \begin{alignat}{2}
      & \text{Indices} & \qquad & \text{Clauses}\nonumber\\
      & i \in \{1, \ldots, 7\} && h_{1, i, 1}
      \label{eq:tape-heads-start}\\
      & i \in \{1, \ldots, 7\},~ j \neq 1 && \neg h_{1, i, j}
      \label{eq:no-more-tape-heads}\\
      & k \in \{1, \ldots, \kappa\} && s_{1, k}^{o_k}
      \label{eq:start-state}\\
      & i \in \{1, \ldots, 7\},~ j \in \{1, \ldots, T\} &&
      \neg c_{1, i, j, 1} \lor \neg c_{1, i, j, 2}
      \label{eq:not-square}\\
      & i \in \{1, \ldots, 5\},~ j < j' \in \{1, \ldots, T\} &&
      \neg c_{1, i, j, 1} \lor c_{1, i, j', 1}
      \label{eq:trailing-empty}\\
      & i \in \{6, 7\},~ j \in \{1, \ldots, T\}
      && c_{1, i, j, 1} \text{ and } \neg c_{1, i, j, 2}
      \label{eq:all-empty}
    \end{alignat}
    This is meant to be read as follows: for each row, the ``Indices'' column
    specifies the range of the indices that the clauses in the ``Clauses'' column
    are quantified over.
    For example, row~\eqref{eq:tape-heads-start} specifies that for all $i \in
    \{1, \ldots, 7\}$, $\varphi_{\mathrm{Boundary}}$ contains the clause $h_{1,
      i, 1}$.
    For the $j \in \{0, T+1\}$ boundary, $\varphi_{\mathrm{Boundary}}$ contains
    the following set of clauses.
    \begin{alignat}{2}
      & \text{Indices} & \qquad & \text{Clauses}\nonumber\\
      & t \in \{1, \ldots, T\},~i \in \{1, \ldots, 7\},
      j \in \{0, T+1\},~k \in \{1, 2\}
      && c_{t, i, j, k} \label{eq:square}
    \end{alignat}
  \end{definition}

  Rows~\eqref{eq:tape-heads-start} and~\eqref{eq:no-more-tape-heads} ensure that
  at time~$t = 1$, each tape has exactly one tape head, and it is located on
  cell $j = 1$.
  Row~\eqref{eq:start-state} ensures that at time~$t=1$, the state is given by
  the start state $\mathsf{start}$.
  (Recall the notation $x_i^o$ to denote the literal $x_i$ if $o = 1$ and $\neg
  x_i$ if $o = 0$.)
  For the remaining rows, we recall that under the encoding of the tape
  alphabet, $\mathrm{enc}_\Gamma(\sqcup) = 10$ and $\mathrm{enc}_\Gamma(\square)
  = 11$.
  As a result, (i) row~\eqref{eq:square} ensures that for all times and tapes,
  the cells $j \in \{0, T+1\}$ contain the~$\square$ symbol, (ii)
  row~\eqref{eq:not-square} ensures that for time~$t=1$, no cell $j \notin \{0,
  T+1\}$ contains the~$\square$ symbol, and (iii) row~\eqref{eq:all-empty}
  ensures that for time~$t=1$ and tapes~$6$ and~$7$, all cells $j \in \{1,
  \ldots, T\}$ contain the $\sqcup$ symbol.
  Finally, row~\eqref{eq:trailing-empty} says that for tapes $i \in \{1, \ldots,
  5\}$, if cell~$j$ contains~$\sqcup$, then every cell $j' > j$ must
  contain~$\sqcup$ as well.
  This means that the five strings encoded by $c_{1, 1}, \ldots, c_{1,5}$ each
  consist of a string of $0$'s and $1$s followed by a string of $\sqcup$'s.
  In short, suppose we write $n$, $x$, $y$, $a$, and~$b$ for the prefixes of
  these strings with no $\sqcup$'s.
  If $\varphi_{\mathrm{Boundary}}$ is satisfied, then the execution tableau
  correctly encodes that the tapes of $\decider$ contain inputs $n$, $x$, $y$,
  $a$, and~$b$ at time~$t=1$.

  Next, we describe the execution constraints.
  These are more complicated than the boundary constraints, and so we will begin
  by describing them in terms of a general Boolean circuit known as the
  \emph{local check circuit}.
  For any time~$t\in \{1, \ldots, T-1\}$ and tape positions $j_1, \ldots, j_7
  \in \{1, \ldots, T\}$, the local check circuit can check that the execution
  tableau properly encodes these tape positions at time~$t+1$ by looking only at
  the encodings of these tape positions and their neighbors (i.e.\ the tape
  positions $j_i \pm 1$ for each~$i \in \{1, \ldots, 7\}$) at time~$t$.

  \begin{definition}\label{def:local-check-circuit}
    In this definition, we will define the \emph{local check circuit}
    $\circuit_{\mathrm{Check}}$.
    It has the following inputs.
    \begin{itemize}
    \item For each $i \in \{1, \ldots, 7\}$, it has the eight inputs
      \begin{equation}
        \label{eq:local-check-inputs}
        \begin{matrix}
          c_{\mathrm{Check}, 0, i, -1} & c_{\mathrm{Check}, 0, i, 0} &
          c_{\mathrm{Check}, 0, i, 1} & \qquad\qquad &
          h_{\mathrm{Check}, 0, i, -1} & h_{\mathrm{Check}, 0, i, 0} &
          h_{\mathrm{Check}, 0, i, 1} \\
          & c_{\mathrm{Check}, 1, i, 0} & & & &
          h_{\mathrm{Check}, 1, i, 0} &
        \end{matrix}
      \end{equation}
      where the $c$-inputs are in $\{0, 1\}^2$ and the $h$-inputs are in $\{0,
      1\}$.
    \item It has two inputs $s_{\mathrm{Check}, 0}, s_{\mathrm{Check}, 1} \in
      \{0, 1\}^\kappa$.
    \end{itemize}
    In addition, for each time $t \in \{1, \ldots, T-1\}$ and tape positions
    $j_1, \ldots, j_7 \in \{1, \ldots, T\}$, we will define a circuit
    $\circuit_{\mathrm{Check}, t, j_1, \ldots, j_7}$ by associating the inputs
    of $\circuit_{\mathrm{Check}}$ with certain variables in~$\mathcal{V}$.
    We do this by associating the following eight inputs from~$\mathcal{V}$ with
    the corresponding variables in Equation~\eqref{eq:local-check-inputs}:
    \begin{equation*}
      \begin{matrix}
        c_{t, i, j_i-1} & c_{t, i, j_i} & c_{t, i, j_i+1} & \qquad\qquad &
        h_{t, i, j_i-1} & h_{t, i, j_i} & h_{t, i, j_i+1} \\
        & c_{t+1, i, j_i} & & & & h_{t+1, i, j_i} &
      \end{matrix}
    \end{equation*}
    as well as by associating $s_t$ and $s_{t+1}$ from~$\mathcal{V}$ with
    $s_{\mathrm{Check}, 0}$ and $s_{\mathrm{Check},1}$, respectively.

    We first define~$\circuit_{\mathrm{Check}}$ as a function. In Proposition~\ref{prop:check-size} below we give an implementation of this function as a circuit. To do so it is convenient to define a function $\circuit_{\mathrm{Check},
      t, j_1, \ldots, j_7}$ for each value of $t, j_1, \ldots, j_7$.
    It will be clear that each of these is in fact the same function
    applied to different inputs, and hence this will
    define~$\circuit_{\mathrm{Check}}$ as well.
    Now, the functions~$\circuit_{\mathrm{Check}, t, j_1, \ldots, j_7}$ perform the following computation. 
    \begin{itemize}
    \item Suppose for each $i \in \{1, \ldots, 7\}$, exactly one of the tape
      positions $j_i-1$, $j_i$, and $j_i+1$ at time~$t$ contains a tape head.
      First, $\circuit_{\mathrm{Check}, t, j_1, \ldots, j_7}$ computes the
      transition function of the Turing machine applied to the contents of
      these~$7$ tape positions, when  the state of~$\decider$ at time~$t$ is $s_{\mathrm{Check}, 0}$. This transition function 
      produces the state of~$\decider$ and contents of these tape positions at
      time~$t+1$, as well as directions to move the~$7$ tape heads in.
      Then $\circuit_{\mathrm{Check}, t, j_1, \ldots, j_7}$ checks that the
      variables for the tape positions $j_1, \ldots, j_7$ and the state
      of~$\decider$ at time~$t+1$ match what they should be according to the computed transition function.
      In addition, if a tape head is marked as moving into a tape cell
      containing a~$\square$ symbol, that tape head remains in place instead.
    \item Suppose for each $i \in \{1, \ldots, 7\}$, none of the tape positions
      $j_i -1$, $j_i$, or $j_i+1$ contain a tape head at time~$t$.
      Then $\circuit_{\mathrm{Check}, t, j_1, \ldots, j_7}$ checks that the
      variables for the tape positions $j_1, \ldots, j_7$ at time~$t+1$ match
      the variables at time $t$.
    \item Otherwise, $\circuit_{\mathrm{Check}, t, j_1, \ldots, j_7}$ accepts.
    \end{itemize}
    This defines $\circuit_{\mathrm{Check}, t, j_1, \ldots, j_7}$, and therefore
    $\circuit_{\mathrm{Check}}$.
  \end{definition}

  Definition~\ref{def:local-check-circuit} shows the utility of introducing the
  boundary variables $c_{t, i, 0}$ and $c_{t, i, T+1}$.
  The function $\circuit_{\mathrm{Check}}$ checks the contents of a cell $c_{t,
    i, j}$ at time $t+1$ by looking at the cell and its neighbors $c_{t, i,
    j-1}, c_{t, i, j+1}$ at time $t$.
  However, those cells with $j \in \{1, T\}$ only have either a left neighbor or
  a right neighbor, and so without the boundary variables we'd have to introduce
  two other local check circuits designed just for these boundary cases.
  The boundary variables then allow us to use the same local check circuit for
  all cells.

  \begin{proposition}\label{prop:check-size}
    A circuit~$\circuit_{\mathrm{Check}}$ of size size at most~$\poly(|\decider|)$ which computes the function described above can be computed in time~$\poly(|\decider|)$.
  \end{proposition}
	
  \begin{proof}
    The circuit has $7 \cdot 12 + 2\cdot \kappa = 84 + 2\kappa$ total Boolean
    inputs, giving a total of $2^{84} \cdot 4^\kappa = O(|K|^2)$ possible input
    strings.
    For each possible fixed input string, the circuit $\circuit_{\mathrm{Check}}$ checks
    if the actual input is equal to the fixed input, which takes $O(\kappa)$
    gates, and then it accepts if the fixed input should be accepting, i.e.\ for each fixed input the correct output is hard-coded in the circuit. Doing so takes $O(|K|^2 \cdot \kappa)$ gates.
    Computing $\circuit_{\mathrm{Check}}$ requires looping over all possible
    input strings and checking which ones are accepting or rejecting.
    This requires computing the transition function of~$\decider$, a task which
    takes time $\poly(|\decider|)$.
    The proposition follows by noting that $|K| \leq |\decider|$.
\end{proof}

\begin{proposition}\label{prop:correct-tableau}
  Suppose that the execution tableau satisfies $\varphi_{\mathrm{Boundary}}$ and
  the circuit $\circuit_{\mathrm{Check}, t, j_1, \ldots, j_7}$, for each $t \in
  \{1, \ldots, T-1\}$ and $j_1, \ldots, j_7 \in \{1, \ldots, T\}$.
  Then the execution tableau correctly encodes the execution of $\decider$ on
  input $(n, x,y, a, b)$ when run for time~$T$.
\end{proposition}
\begin{proof}
  To show this, we will show for each time $t \in \{1, \ldots T\}$ that the
  variables in the execution tableau corresponding to time~$t$ correctly encode
  the state of~$\decider$ and the contents of the seven tapes at time~$t$.
  The proof is by induction on~$t$.
  The base case of~$t=1$ follows from the tableau satisfying
  $\varphi_{\mathrm{Boundary}}$.

  Next we perform the induction step.
  Assuming the statement holds for time $t \in \{1, \ldots, T-1\}$, we will show
  it holds for time~$t+1$ as well.
  Let $i^* \in \{1, \ldots, 7\}$ be a tape, and consider a tape position
  $j_{i^*} \in \{1, \ldots, T\}$.
  We will show that the variables correctly encode the contents of this tape
  position at time $t+1$.
  Suppose one of the tape positions $j_{i^*}-1$, $j_{i^*}$, or $j_{i^*}+1$ at
  time $t$ has a tape head.
  For each of the other tapes $i \neq i^*$, we select a tape position $j_{i}$
  such that either $j_{i}-1$, $j_{i}$, or $j_{i}+1$ has a tape head at time $t$.
  (These positions are guaranteed to exist since each tape has exactly one tape
  head.)
  
  By the induction hypothesis, for each tape $i \in \{1, \ldots, 7\}$ the
  variables corresponding to tape cells $j_i-1$, $j_i$, and $j_i+1$ correctly
  encode the contents of these cells at time~$t$.
  By assumption, $\circuit_{\mathrm{Check}, t, j_1, \ldots, j_7}$ evaluates
  to~$1$.
  In this case, it calculates the transition function of~$\decider$ to compute
  the contents of the tape cells $j_1, \ldots, j_7$ at time~$t+1$ and checks
  that the corresponding variables encode these contents.
  As a result, the tape position $j_{i^*}$ on tape~$i^*$ is correctly encoded.
  In addition, it computes the state of~$\decider$ at time~$t+1$ and checks that
  the corresponding variables encode this state.
  This completes the induction step.
  The case when none of the tape positions $j_{i^*}-1$, $j_{i^*}$, and
  $j_{i^*}+1$ at time $t$ contain a tape head follows similarly.
  Finally, the variables for all tape positions $j \in \{0, T+1\}$ are correctly
  encoded due to $\varphi_{\mathrm{Boundary}}$ being satisfied.
\end{proof}
  
Proposition~\ref{prop:correct-tableau} gives a set of constraints that ensure
the execution tableau properly encodes the execution of~$\decider$.
Our next step will be to convert these constraints into a single 3SAT formula.
This entails transforming each circuit $\circuit_{\mathrm{Check}, t, i_1,
  \ldots, i_7}$ into a 3SAT formula.
We do so using the following reduction.

\begin{proposition}[Circuit-to-3SAT]\label{prop:circuit-to-sat-reduction}
  There is an algorithm which, on input a size-$r$ circuit~$\circuit$ on
  variables $x \in \{0, 1\}^n$, runs in time $\poly(r)$ and outputs a 3SAT
  formula $\varphi$ on variables $x \in \{0, 1\}^n$ and $y \in \{0, 1\}^r$ with
  $O(r)$ clauses such that for all $x$, $\circuit(x) = 1$ if and only if there
  exists a $y$ such that $x$ and~$y$ satisfy $\varphi$.
\end{proposition}
\begin{proof}
  This is the textbook circuit-to-3SAT reduction.
  For each gate $i \in \{1, \ldots, r\}$, the algorithm introduces a variable
  $y_i \in \{0, 1\}$, so that the total collection of variables is
  $\mathcal{V}_{\mathrm{total}} = \{x_i\}_{i \in \{1, \ldots, n\}} \cup
  \{y_i\}_{i \in \{1, \ldots, r\}}$.
  Consider gate~$i \in \{1, \ldots, r\}$, and let $z_1, z_2 \in
  \mathcal{V}_{\mathrm{total}}$ be the pair of variables feeding into it.
  If gate~$i$ is an AND gate, then~$\varphi$ includes the constraints
  \begin{equation}\label{eq:AND-gate}
    (\neg z_1 \lor \neg z_2 \lor y_i)\;,
    \quad
    (z_1 \lor  z_2 \lor \neg y_i)\;,
    \quad
    (z_1 \lor  \neg z_2 \lor \neg y_i)\;,
    \quad
    (\neg z_1 \lor  z_2 \lor \neg y_i)\;.
  \end{equation}
  These constraints are satisfied if and only if $y_i = z_1 \land z_2$.
  The case of gate~$i$ being an OR gate follows similarly.
  Finally, $\varphi$ includes the constraint $(y_{i^*})$, where $i^* \in \{1,
  \ldots, r\}$ is the output gate.
  It follows that $x, y$ satisfy $\varphi$ if and only if $\circuit(x) = 1$ and
  for each $i \in \{1, \ldots, r\}$, $y_i$ is the value computed by gate~$i$ in
  circuit~$\circuit$ on input~$x$.
  In total, $\varphi$ has $4r+1$ clauses and is computable in time $\poly(r)$.
\end{proof}

We now apply the algorithm from Proposition~\ref{prop:circuit-to-sat-reduction}
to $\circuit_{\mathrm{Check}}$, which by Proposition~\ref{prop:check-size} has
size $r \leq \poly(|\decider|)$.
This produces a 3SAT formula $\varphi_{\mathrm{Check}}$ on the variables in
$\circuit_{\mathrm{Check}}$ plus auxiliary variables $a_k$ for $k \in \{1,
\ldots, r\}$ added by the reduction.
Now, for each $t \in \{1, \ldots, T-1\}$ and $j_1, \ldots, j_7 \in \{1, \ldots,
T\}$, we define a 3SAT formula $\varphi_{\mathrm{Check}, t, j_1, \ldots, j_7}$
analogously to $\circuit_{\mathrm{Check}, t, j_1, \ldots, j_7}$.
We begin by associating those variables in $\varphi_{\mathrm{Check}}$ which come
from $\circuit_{\mathrm{Check}}$'s inputs with the variables in~$\mathcal{V}$ as
in Definition~\ref{def:local-check-circuit}.
Next, for each $k \in \{1, \ldots, r\}$, we introduce a new variable $a_{t, j_1,
  \ldots, j_7, k} \in \{0, 1\}$ and associate it with the variable $a_k$.
This defines $\varphi_{\mathrm{Check}, t, j_1, \ldots, j_7}$.
In summary, the final 3SAT instance produced by the Cook-Levin reduction is
\begin{equation}\label{eq:cook-levin-3sat}
  \varphi:= \varphi_{\mathrm{Boundary}} \land \big(\bigwedge_{t, j_1, \ldots, j_7}
  \varphi_{\mathrm{Check}, t, j_1, \ldots, j_7}\big)\;.
\end{equation}
By Proposition~\ref{prop:circuit-to-sat-reduction}, the execution tableau
properly encodes the execution of~$\decider$ if and only if there exists a
setting to the auxiliary variables satisfying the 3SAT formula $\varphi$.
In total, $\varphi$ contains
\begin{equation}\label{eq:variable-count}
|\mathcal{V}| + O(T^8) \cdot r
\leq O\big(T^2 + T \cdot \log(|\decider|)
    + T^8\cdot \poly(|\decider|)\big)
 = \poly(T, |\decider|)
\end{equation}
variables.
The second term on the left-hand side of Equation~\eqref{eq:variable-count}
corresponds to the auxiliary variables in each of the $O(T^8)$ copies of
$\varphi_{\mathrm{Check}}$.
The inequality follows by Equation~\eqref{eq:tableau-variable-count} and the
fact that $r \leq \poly(|\decider|)$.
  
Our next task is to represent the 3SAT formula~$\varphi$ succinctly.
To do this, we will provide a circuit~$\circuit$ which succinctly describes a
3SAT formula~$\varphi_{\mathrm{3SAT}}$ which, while not literally equal
to~$\varphi$, will be isomorphic to it.
This means that, for example, $\varphi_{\mathrm{3SAT}}$ may not even have the
same number of variables as~$\varphi$, but any variable in~$\varphi$ will
correspond in a clear and direct manner to a variable
in~$\varphi_{\mathrm{3SAT}}$, and any remaining variables
in~$\varphi_{\mathrm{3SAT}}$ do not appear in any clauses.
This circuit is constructed as follows.

\begin{definition}
  In this definition we construct the circuit~$\circuit$.
  It has three inputs~$z_1, z_2, z_3$ of length~$m$, which we specify below in
  Equation~\eqref{eq:length-of-m}, and three inputs~$o_1, o_2, o_3 \in \{0,
  1\}$.
	For each~$\nu \in \{1, 2, 3\}$, each $z_\nu$ is supposed to specify a variable
  in $\varphi$ according to a format we will now specify.
	If~$z_\nu$ is not properly formatted, then it does not correspond to a
  variable in~$\varphi$; if any of~$z_1, z_2, z_3$ is not properly formatted,
  then~$\circuit$ automatically outputs~$0$.
	Below, we will often write substrings of the $z_\nu$'s as though they are
  integers from some specified range, i.e.\ $a \in \{b, \ldots, c\}$.
	This means that $a$ is represented as a binary string of length $\lceil
  \log(c+1) \rceil$, which is to be interpreted as the binary encoding of an
  integer between~$b$ and~$c$.
	
  The input $z_\nu$ is formatted as a string $(\omega, \alpha, \beta_1, \beta_2,
  \beta_3, \beta_4)$.
	The first substring $\omega$ has length $m - (|\alpha| + |\beta_1| + \cdots +
  |\beta_4|)$ bits and is formatted to be the all-zeroes string.
	Its purpose is to pad the inputs to have the length~$m$ we specify below.
	Next, $\alpha$ is formatted as an integer $\alpha \in \{1,2, 3, 4\}$.
	The variable encoded by~$z_\nu$ is specified by $\beta_\alpha$, and the other
  three $\beta$'s should be the all-zeros string.
	We now specify the encoding of $\beta_\alpha$, conditioned on the value
  of~$\alpha$.
	\begin{enumerate}
	\item $\beta_1$ is formatted as $(t, i, j, k)$, where\label{item:a=1}
		\begin{equation*}
      t \in \{1, \ldots, T\}, \quad i \in \{1, \ldots, 7\},
      \quad j \in \{0, 1, \ldots, T+1\}, \quad k \in \{1, 2\}.
		\end{equation*}
		This corresponds to the variable $c_{t, i, j, k}$.
	\item $\beta_2$ is formatted as $(t, i, j)$, where\label{item:a=2}
		\begin{equation*}
      t \in \{1, \ldots, T\}, \quad i \in \{1, \ldots, 7\},
      \quad j \in \{0, 1, \ldots, T+1\}.
		\end{equation*}
		This corresponds to the variable $h_{t, i, j}$.
	\item $\beta_3$ is formatted as $(t, k)$, where\label{item:a=3}
		\begin{equation*}
      t \in \{1, \ldots, T\}, \quad k \in \{1, \ldots, \kappa\}.
		\end{equation*}
		This corresponds to the variable $s_{t, k}$.
	\item $\beta_4$ is formatted as $(t, j_1, \ldots, j_7, k)$, where\label{item:a=4}
		\begin{equation*}
      t \in \{1, \ldots, T-1\}, \quad j_1, \ldots, j_7 \in \{1, 2, \ldots, T\},
      \quad k \in \{1, \ldots, r\}.
		\end{equation*}
		This corresponds to the variable $a_{t, j_1, \ldots, j_7, k}$.
	\end{enumerate}
	In total, the length of these substrings is
	\begin{equation*}
    |\alpha| + |\beta_1| + \cdots + |\beta_4| = O(\log(T) + \log(\kappa) +
    \log(r)) = O(\log(T) + \log(|\decider|))\;.
	\end{equation*}
	As a result, because $|\decider| \leq \sigma$, the length of the
  ``padding''~$\omega$ can be chosen so that each~$z_i$ has length
	\begin{equation}\label{eq:length-of-m}
    m = O(\log(T) + \log(\sigma))\;.
	\end{equation}
	Checking that each $z_\nu$ is properly formatted can be done by checking that
  certain substrings of $z_1$, $z_2$, and $z_3$ encode integers which fall
  within specified ranges.
	This can be done using $O(m)$ gates.

  Having specified the inputs, we can now specify the execution of the circuit,
  and we may assume that $z_1, z_2, z_3$ are properly formatted.
  Implementing the clauses from $\varphi_{\mathrm{Boundary}}$ is simple; we
  specify how to implement the clause from Equation~\eqref{eq:tape-heads-start}.
  \begin{itemize}
  \item Suppose for input $z_1$, $\alpha = 2$.
    Then $\beta_2$ can be parsed as $(t, i, j)$.
    The circuit accepts if $t=j = 1$ and $o_1 = 1$, regardless of~$z_2$ and~$z_3$.
    This ensures that $\varphi_{\mathrm{3SAT}}$ includes $x_{z_1} \lor
    x_{z_2}^{o_2} \lor x_{z_3}^{o_3}$ for any $z_2, z_3, o_2, o_3$, which is
    equivalent to including the arity-one clause $x_{z_1}$.
  \end{itemize}
  This can be implemented with $O(m)$ gates.
  Similar arguments can be used to implement the clauses from
  Equations~\eqref{eq:no-more-tape-heads}-\eqref{eq:square} using $O(m)$ gates
  apiece.

  Implementing the clauses from the formulas $\varphi_{\mathrm{Check}, t, j_1,
    \ldots, j_7}$ is more challenging.
  From Equation~\eqref{eq:AND-gate}, we can see that any constraint in this
  formula always involves a variable of the form $a_{t, j_1, \ldots, j_7, k}$
  for some $k \in \{1, \ldots, r\}$.
  \begin{enumerate}
  \item First check if one of its inputs $z_1, z_2, z_3$ corresponds to such a
    variable.
    This can be done with $O(1)$ gates simply by checking if for any of the
    $z_i$'s, $a = 4$.
    If so, this specifies the values of $t, j_1, \ldots, j_7$.
  \item Each variable in $\varphi_{\mathrm{Check}, t, j_1, \ldots, j_7}$ is
    associated with a variable in $\varphi_{\mathrm{Check}}$.
    The circuit checks if all the $z_i$'s are contained in
    $\varphi_{\mathrm{Check}, t, j_1, \ldots, j_7}$ and then it computes which
    variable in $\varphi_{\mathrm{Check}}$ they are associated with.
    We include below the example of checking whether~$z_1$ is associated with
    the variable $c_{\mathrm{Check},1, 1, 0, 0}$.
    \begin{itemize}
    \item The circuit~$\circuit$ first computes $t+1$, which takes $O(m)$ gates.
      It then looks at~$z_1$ and checks if~$\alpha=1$.
      If so, then $\beta_1 = (t', i', j', k')$, and so it tests the equalities
      $t' = t+1$, $i' = 1$, $j' = j_1$, and $k' = 0$, each of which takes $O(m)$
      gates to test.
      If so, then~$z_1$ is associated with $c_{\mathrm{Check},1, 1, 0, 0}$.
    \end{itemize}
    There are $\poly(|\decider|)$ variables in $\varphi_{\mathrm{Check}}$ and,
    for each variable $z_i$, it takes $O(m)$ gates to determine whether $z_i$ is
    associated with this variable.
    As a result, because $|\decider| \leq \sigma$, this takes $\poly(\sigma,
    \log(T))$ gates to compute.
  \item Whether $z_1$, $z_2$, and $z_3$ share a clause in
    $\varphi_{\mathrm{Check}, t, j_1, \ldots, j_7}$ depends only on which
    variables in $\varphi_{\mathrm{Check}}$ they are associated with.
    As a result, after having computed these variables, the algorithm can
    hard-code whether the circuit~$\circuit$ should accept.
  \end{enumerate}
  This completes the description of~$\circuit$.
  In total, it contains $\poly(\sigma, \log(T))$ gates.
  Computing $\circuit$ first requires computing $\varphi_{\mathrm{Check}}$,
  which takes time $\poly(|\decider|) \leq \poly(\sigma)$ by
  Propositions~\ref{prop:check-size} and~\ref{prop:circuit-to-sat-reduction}.
  After that, the steps outlined above for construction~$\circuit$ take
  time~$\poly(\sigma, \log(T))$.
\end{definition}

The circuit~$\circuit$ succinctly describes~$\varphi$ in the loose sense
described above.
Recall that~$\varphi$ accepts if and only if it encodes the execution
of~$\decider$ up to time~$T$.
We now modify~$\circuit$ to (i) hard code~$n$, $x$, and~$y$ onto~$\decider$'s
input tapes, and (ii) ensure that~$\decider$ accepts.
We demonstrate how to do so by hard coding~$n$ as an example.
\begin{itemize}
\item The input~$n$ is described by some string $\nu$ of length~$\ell =
  O(\log(n))$.
	We would like to hard-code the string $(\nu, \sqcup^{T - \ell})$ into the
  first tape of~$\decider$.
	To do this, we first check if~$z_1$ corresponds to the variable $c_{t, i, j,
    k}$ for some values of~$t, i, j, k$.
	Then, we check if $t = 1$, the time where the inputs appear on the tapes, and
  $i = 1$, corresponding to the first tape.
	If so, we branch on whether $j \leq \ell$.
	If it is, then if $k = 1$ the circuit accepts if $o_1 = 0$, and if $k = 2$ the
  circuit accepts if $o_1 = \nu_j$.
	Otherwise, if $j > \ell$, then if $k = 1$ the circuit accepts if $o_1 = 1$,
  and if $k = 2$ the circuit accepts if $o_1 = 0$.
	This can be done with $\poly(\log(n), \log(T))$ gates.
\end{itemize}
This modifies~$\varphi$ so that it only accepts if~$n$ is written on its first
tape at input.
We can similarly hard-code~$x$ and~$y$ onto the second and third input tapes and
hard-code the accepting state as the final state of~$\decider$.
In total, this takes $\poly(\log(n), |x|, |y|, \log(T), \sigma)$, which is
$\poly(\log(n), Q, \log(T), \sigma)$ because~$x$ and~$y$ have length at
most~$Q$.

It remains to ensure that the variables corresponding to the~$4$th and~$5$th at
time~$t = 1$ are the lexicographically-first named variables in~$\varphi$.
However, this is simple and can be done using $\poly(\log(n), Q, \log(T),
\sigma)$ gates.
This concludes the construction.
\end{proof}

\subsection{A succinct 5SAT description for deciders}
\label{sec:succinct-deciders}

Proposition~\ref{prop:standard-succinct-sat} allows us to convert any decider
$\decider$ and inputs $n, x, y$ into a 3SAT formula $\varphi_{\mathrm{3SAT}}$,
succinctly described by a circuit~$\circuit$, which represents it.
However, there are two undesirable properties of this construction, which we
describe below.
\begin{enumerate}
\item First, evaluating any clause $(w_{i_1}^{o_{i_1}} \lor
  w_{i_2}^{o_{i_2}}\lor w_{i_3}^{o_{i_3}})$ of $\varphi_{\mathrm{3SAT}}$
  requires evaluating the \emph{same} assignment~$w$ at three separate points.
	While this is fine when the assignment~$w$ is provided in full to the
  verifier, it can be a problem when the verifier is only able to query the
  points in~$w$ by interacting with a prover.
	In this case, the verifier might send the prover the values $i_1, i_2, i_3$,
  who responds with three bits $b_1, b_2, b_3 \in \{0, 1\}$, purported to be the
  values $w_{i_1}, w_{i_2}, w_{i_3}$ for some assignment $w \in \{0, 1\}^M$.
	As it will turn out, in the answer reduced protocol below, the verifier will
  actually be able to force the prover to reply using \emph{three} different
  assignments.
	In other words, the prover will have three assignments $w_1, w_2, w_3 \in \{0,
  1\}^M$ such that, for any $i_1, i_2, i_3$ provided to it by the verifier, it
  will respond with $b_1 = w_{1, i_1}, b_2 = w_{2, i_2}, b_3 = w_{3, i_3}$.
	However, even given this there is no guarantee that the three assignments are
  the \emph{same} assignments (i.e.\ that $w_1 = w_2 = w_3$).
	In the work of~\cite{NW19}, this was accomplished by an additional subroutine
  called the \emph{intersecting lines test}, which would enforce consistency
  between $w_1$, $w_2$, and $w_3$.
	In this work, on the hand, we would like to relax the assumption that $w_1$,
  $w_2$, and $w_3$ must be the same.
	This will allow us to not use the intersecting lines test, simplifying the
  answer reduction protocol.
	(In fact, the answer reduced verifier will not be querying the assignments
  directly, but rather low-degree encodings of these assignments.)
\item Second, we are guaranteed that $w = (a, b, c)$ satisfies
  $\varphi_{\mathrm{3SAT}}$ if and only if $\decider$ accepts $(n, x, y, a, b)$.
	However, it is inconvenient that~$a$ and~$b$ are contained as substrings
  of~$w$.
	To see why, recall from Section~\ref{sec:oracle} that the oracularized
  verifier sometimes gives one prover a pair of questions $(x, y)$ and another
  prover just one of the questions---say, $x$.
	The answer reduced verifier will sample its questions similarly; as for its
  answers, it might expect the first prover to respond with a string $w = (a, b,
  c)$ that satisfies $\varphi_{\mathrm{3SAT}}$ and the second prover to respond
  with a string $a'$ such that $a = a'$.
	Verifying that $a = a'$ requires the verifier to sample a uniformly random
  point from~$w$, restricted to the coordinates in~$a$.
	As it turns out, generating a uniform point from a substring is extremely
  cumbersome, though not impossible, to do when the verifier's questions are
  expected to be sampled from conditional linear functions.
	To remove this complication, though, we would instead prefer if the provers'
  answers were formatted in a way that gave the verifier direct access to the
  strings~$a$ and~$b$.
\end{enumerate}
  
In the remainder of this section, we will show how to modify the Succinct-3SAT
circuits produced by Proposition~\ref{prop:standard-succinct-sat} in order to
ameliorate these two difficulties.
Doing so entails modifying the $\varphi_{\mathrm{3SAT}}$ formula from
Proposition~\ref{prop:standard-succinct-sat} to produce a \emph{5SAT formula}
$\varphi_{\mathrm{5SAT}}$.
The clauses of $\varphi_{\mathrm{5SAT}}$ will be of the form
\begin{equation*}
  a_{i_1}^{o_1} \lor b_{i_2}^{o_2} \lor w_{1, i_3}^{o_3}
  \lor w_{2, i_4}^{o_4} \lor w_{3, i_5}^{o_5}\;,
\end{equation*}
where $a$, $b$, $w_1$, $w_2$, and $w_3$ are five separate assignments which are
not assumed to be equal.
The guarantee is that $(a, b, w_1, w_2, w_3)$ satisfies
$\varphi_{\mathrm{5SAT}}$ if and only if $\decider$ accepts $(n, x, y, a, b)$.
This addresses the two concerns from above: each clause is totally
\emph{decoupled}, meaning it samples~$5$ variables from~$5$ different
assignments, and so no consistency check must be performed between the
assignments.
In addition, the first two strings exactly correspond to~$a$ and~$b$, addressing
the second item.

We now formally define decoupled 5SAT instances and how they succinctly
represent bounded deciders.
Following that, we show how the succinct 3SAT instance $\varphi_{\mathrm{3SAT}}$
produced by Proposition~\ref{prop:standard-succinct-sat} can be modified to
produce a succinct 5SAT instance $\varphi_{\mathrm{5SAT}}$ which represents
$\decider$.

\begin{definition}[Decoupled 5SAT and its succinct descriptions]
  A \emph{block} of variables~$x_i$ is a tuple $x_i = (x_{i,0}, \ldots, x_{i,
    N_i-1})$.
  A formula $\varphi$ on~$5$ blocks $x_1, x_2, \ldots, x_5$ of variables is
  called a \emph{decoupled 5SAT formula} if every clause is of the form
  \begin{equation}\label{eq:5sat}
    x_{1,\, i_1}^{o_1}
    \lor x_{2,\, i_2}^{o_2}
    \lor x_{3,\, i_3}^{o_3}
    \lor x_{4,\, i_4}^{o_4}
    \lor x_{5,\, i_5}^{o_5}\;,
  \end{equation}
  for $i_j \in \{0, 1, \ldots, N_j-1\}$ and $o_1, \ldots, o_5 \in \{0, 1\}$.
  (Recall from Definition~\ref{def:succinct-formulas} that the notation $x^{o}$
  means $x$ if $o = 1$ and $\neg x$ if $o = 0$.)

  For each $i\in\{1, 2, \ldots, 5\}$, suppose each $N_i$ is a power of two, and write
  it as $N_i = 2^{n_i}$.
  Let $\circuit$ be a circuit with five inputs of length $n_1, n_2, \ldots, n_5$ and
  five single-bit inputs.
  Then $\circuit$ \emph{succinctly describes decoupled $\varphi$} if, for all
  $i_j \in \{0, 1, \ldots, N_j-1\}$ and $o_1, o_2, \ldots, o_5 \in \{0, 1\}$,
  \begin{equation}\label{eq:put-in-numbers}
    \circuit(i_1, i_2, \ldots, i_5,
    o_1, o_2, \ldots, o_5) = 1
  \end{equation}
  if and only if the clause in~\eqref{eq:5sat} is included in~$\varphi$.
  As in Definition~\ref{def:succinct-formulas}, we slightly abuse notation and
  use the convention that a number~$a$ between~$0$ and~$2^{n_i}-1$ is
  interpreted as its binary encoding $\binary{n_i}(a)$ when provided as input to
  a set of $n_i$ single-bit wires.
\end{definition}

\begin{definition}[Succinct descriptions for bounded deciders]
  Let $\decider$ be a decider.
  Fix an index $n \in \N$ and a time $T \in \N$.
  Let $L = 2^\ell$ be a power of two that is at least as large as~$2T$.
  Let~$x$ and~$y$ be strings, $r\in\N$ and $R=2^r$.

  Consider a circuit~$\circuit$ with two inputs of length~$\ell$, three inputs
  of length~$r$, and $5$ single-bit inputs.
  Let $\varphi_C$ be the decoupled 5SAT instance with two blocks of variables of
  size~$L$ and three blocks of size~$R$ which~$\circuit$ succinctly describes.
  Then we say that $\circuit$ \emph{succinctly describes $\decider$ (on
    inputs~$n$, $x$, and~$y$ and time~$T$)} if, for all $a, b \in \{0, 1\}^{L}$,
  there exists $w_1, w_2, w_3 \in \{0, 1\}^{R}$ such that $a, b, w_1, w_2, w_3$
  satisfy $\varphi_{\circuit}$ if and only if there exist $a_{\mathrm{prefix}},
  b_{\mathrm{prefix}} \in \{0, 1\}^*$ of lengths~$\ell_a, \ell_b \leq T$,
  respectively, such that
  \begin{equation*}
    a = \mathrm{enc}_\Gamma(a_{\mathrm{prefix}}\,,\, \sqcup^{L/2 - \ell_a})
  \quad
  \text{and}
  \quad
  b = \mathrm{enc}_\Gamma(b_{\mathrm{prefix}}\,,\, \sqcup^{L/2 - \ell_b})
  \end{equation*}
  and $\decider$ accepts $(n, x, y, a_{\mathrm{prefix}}, b_{\mathrm{prefix}})$
  in time~$T$.
\end{definition}

In this definition of succinct descriptions, the answers~$a$ and~$b$ are
isolated, in that the first input of~$\circuit$ of length~$\ell$ indexes
into~$a$ and the second input of length~$\ell$ indexes into~$b$.
The next proposition shows how to construct such descriptions.

\begin{proposition}[Explicit succinct descriptions]
  \label{prop:explicit-succinct-deciders}
  There is a Turing machine $\succinctdecider$ with the following properties.
  Let $\decider$ be a decider, let $n$, $T$,~$\qlen$, and $\sigma$ be integers
  with $\qlen \leq T$ and $|\decider| \leq \sigma$, and let~$x$ and~$y$ be
  strings of length at most~$\qlen$.
  Then on input $(\decider, n, T, \qlen, \sigma, x, y)$, $\succinctdecider$
  outputs a circuit~$\circuit$ with two inputs of length~$\ell_0(T)$, three of
  length $r_0(T, \sigma)$, and five single-bit inputs which succinctly
  describes $\decider$ on inputs~$n$, $x$, and~$y$ and time~$T$.
  Moreover, the following hold.
  \begin{enumerate}
  \item $\ell_0(T) = \lceil \log(2T) \rceil$.
  \item $r_0(T,\sigma) = O(\log(T)+\log(\sigma))$,
  \item $\circuit$ has at most $s_0(n,T,\qlen,\sigma) = \poly(\log(T), \log(n),
    \qlen,\sigma)$ gates,
  \item $\succinctdecider$ runs in time $\poly(\log(T), \log(n),
    \qlen,\sigma)$, and the parameters $\ell_0, r_0, s_0$ can be computed
    from $n, T, \qlen,\sigma$ in time $\poly(\log(T), \log(n),
    \log(\qlen),\sigma)$.
  \end{enumerate}
\end{proposition}

\begin{proof}
  The Turing machine $\succinctdecider$ begins by running the algorithm in
  Proposition~\ref{prop:standard-succinct-sat} on input
  $(\decider,n,T,\qlen,\sigma, x,y)$ to produce a
  circuit~$\circuit_{\mathrm{3SAT}}$ on~$3r_0+3$ inputs, where $r_0 =
  O(\log(T)+\log(\sigma))$ is the parameter $m$ from the proposition.
  Set $\ell_0 = \lceil \log(2T) \rceil$, $L = 2^{\ell_0}$ and $R = 2^{r_0}$.
  Given this, $\succinctdecider$ returns the circuit~$\circuit$ with inputs
  $i_1, i_2 \in \{0, 1, \ldots, L-1\}$, $i_3, i_4, i_5 \in \{0, 1, \ldots,
  R-1\}$, $o_1, o_2, \ldots, o_5 \in \{0, 1\}$, and
  \[
    \circuit(i_1, i_2, \ldots, i_5, o_1, o_2, \ldots, o_5) = 1,
  \]
  if one of the following conditions hold.
  \begin{align*}
    & \circuit_{\mathrm{3SAT}}(i_3, i_4, i_5, o_3, o_4, o_5) = 1\;,\\
    & (i_1 < 2T) \land (i_1 = i_3) \land (o_1 \neq o_3)\;,\\
    & (i_2 < 2T)\land (i_2 = i_3 - 2T) \land (o_2 \neq o_3)\;,\\
    & (i_1 \geq 2T) \land (\text{$i_1$ is odd}) \land (o_1 = 1)\;,\\
    & (i_1 \geq 2T) \land (\text{$i_1$ is even}) \land (o_1 = 0)\;,\\
    & (i_2 \geq 2T) \land (\text{$i_2$ is odd}) \land (o_2 = 1)\;,\\
    & (i_2 \geq 2T) \land (\text{$i_2$ is even}) \land (o_2 = 0)\;,\\
    & (i_3 = i_4) \land (o_3 \neq o_4)\;,\\
    & (i_4 = i_5) \land (o_4 \neq o_5)\;.
  \end{align*}
  It is not hard to verify that testing ``$(i_1 < 2T)$'' can be done with
  $O(\ell_0)$ AND and OR gates, and testing $(i_2 = i_3 - 2T)$ can be done with
  $O(m)$ AND and OR gates.
  Using similar estimates for the remaining sub-circuits, we compute
  \begin{equation*}
    \begin{split}
      \mathrm{size}(\circuit) = \mathrm{size}(\circuit_{\mathrm{3SAT}}) +
      O(r_0 + \ell_0) & = \mathrm{size}(\circuit_{\mathrm{3SAT}}) + O(r_0) \\
      & \leq \poly(\log(T), \log(n), \qlen,\sigma) + O(\log(T) + \log(\sigma))\;.
    \end{split}
  \end{equation*}
  In addition, due to the simplicity of these modifications, we conclude that
  the runtime of $\succinctdecider$ is dominated by the runtime of the algorithm
  from Proposition~\ref{prop:standard-succinct-sat}, which is $\poly(\log(T),
  \log(n), \qlen,\sigma)$.

  Now we show that $\circuit$ succinctly describes $\decider$ on inputs~$n$,
  $x$, and~$y$ and time~$T$.
  To begin, we describe the decoupled 5SAT formula $\varphi_{\circuit}$.
  Let us first consider the constraints in $\varphi_{\circuit}$ which are
  implied by the final constraint, i.e.\ those of the form
  \begin{equation*}
    a_{i_1}^{o_1} \lor b_{i_2}^{o_2} \lor (w_1)_{i_3}^{o_3} \lor (w_2)_{i_4}^{o_4}
    \lor (w_3)_{i_5}^{o_5}
  \end{equation*}
  whenever $i_4 = i_5$ and $o_4 \neq o_5$.
  For any fixed $i_1, i_2, i_3$, the negations $o_1, o_2, o_3$ can take any
  values, and as a result, the first three bits in the constraint vary over all
  assignments in $\{0, 1\}^3$.
  This means that these constraints are satisfied if and only if
  $(w_2)_{i_4}^{o_4} \lor (w_3)_{i_5}^{o_5}$ is satisfied whenever $i_4 = i_5$
  and $o_4 \neq o_5$.
  This, in turn, is equivalent to the constraint $w_2 = w_3$.
  Carrying out similar arguments for the entire circuit, we can express the
  formula~$\varphi_{\circuit}$ as follows.
  \begin{multline*}
    \varphi_{\circuit}(a, b, w_1, w_2, w_3)
    =
    \varphi_{\mathrm{3SAT}}(w_1, w_2, w_3)
    \land (w_{1, 1} = a_1) \land (w_{1, 2} = b_1)\\
    \land (a_2= (10)^{L/2 -T}) \land (b_2 = (10)^{L/2-T})
    \land (w_1 = w_2)
    \land (w_2 = w_3).
  \end{multline*}
  Here, we write $\varphi_{\mathrm{3SAT}}(w_1, w_2, w_3)$ for the formula in
  which, for each constraint in $\varphi_{\mathrm{3SAT}}$, the first variable is
  taken from $w_1$, the second from $w_2$, and the third from $w_3$.
  In addition, we write $a = (a_1, a_2)$, where $a_1$ is the first~$2T$ bits
  in~$a$ and $a_2$ is the remaining $L-2T$ bits, and similarly for~$b = (b_1,
  b_2)$.
  We also write $w_1 = (w_{1,1}, w_{1,2}, w_{1,3})$, where $w_{1, 1}$ contains
  the first~$2T$ bits in~$w_1$, $w_{1,2}$ contains the second~$2T$ bits,
  and~$w_{1,3}$ contains the remaining $R-4T$ bits.
  As a result, $\varphi_{\circuit}$ is satisfied only if $w_1 = w_2 = w_3 =
  (a_1, b_1, c)$ for some string $c \in \{0, 1\}^{R-4T}$.
  In this case, calling $w = (a_1, b_1, c)$, $\varphi_{\circuit}$ is satisfied
  only if $\varphi_{\mathrm{3SAT}}(w)$ is.
  By Proposition~\ref{prop:standard-succinct-sat}, this implies that there exists a
  string~$a_{\mathrm{prefix}}$ of length $\ell_a \leq T$ such that $ a_1 =
  \mathrm{enc}_\Gamma(a_{\mathrm{prefix}},\,\sqcup^{T-\ell_a})$.
  This, in turn, implies that
  \begin{equation*}
    a = (a_1, a_2) = (\mathrm{enc}_\Gamma(a_{\mathrm{prefix}},\,
    \sqcup^{T-\ell_a}), (10)^{L-2T})
    = \mathrm{enc}_\Gamma(a_{\mathrm{prefix}},\,\sqcup^{L/2-\ell_a})\;,
  \end{equation*}
  using the fact that $\mathrm{enc}_\Gamma(\sqcup) = 10$, and similarly for~$b$.
  Finally, Proposition~\ref{prop:standard-succinct-sat} implies that $\decider$
  accepts $(n, x, y, a_{\mathrm{prefix}}, b_{\mathrm{prefix}})$ in time~$T$.
  This completes the proof.
\end{proof}

As stated above, moving from 3SAT to 5SAT allows us to devote the first two
inputs to~$a$ and~$b$.
In addition, we have added extra constraints into $\varphi_{\circuit}$ which
enforce that $w_1 = w_2 = w_3$, which means that we can relax this assumption on
these assignments.

We now show a simple transformation that takes in a succinct circuit~$\circuit$
and outputs another succinct circuit~$\circuit'$ whose~$5$ inputs are ``padded''
to contain more input bits.




\begin{proposition}[Padding]\label{prop:exciting-padding-prop}
  Let $\circuit$ be a circuit of size~$s$ with two inputs of length~$\ell$,
  three inputs of length~$r$, and~$5$ single-bit inputs.
  Suppose $\circuit$ succinctly describes $\decider$ on inputs $n$, $x$, and $y$
  and time~$T$.
  Then there is an algorithm which takes as input $(\circuit, \ell, r, \ell',
  r')$, with $\ell' \geq \ell$ and $r' \geq r$, and in time~$\poly(s, \ell',
  r')$ outputs a circuit~$\circuit'$ with the following properties.
  First, $\circuit'$ has two inputs of length~$\ell'$, three inputs of
  length~$r'$, and~$5$ single-bit inputs, and its size is $s + \poly(\ell',
  r')$.
  Second, it succinctly describes $\decider$ on inputs~$n$, $x$, and~$y$ and
  time~$T$.
\end{proposition}
\begin{proof}
  Write $L = 2^\ell, R = 2^r$ and $L' = 2^{\ell'}, R' = 2^{r'}$.
  The algorithm constructs the circuit~$\circuit'$ which on inputs $i_1, i_2 \in
  \{0, \ldots, L'-1\}$, $i_3, i_4, i_5 \in \{0, \ldots, R'-1\}$, and $o_1,
  \ldots, o_5 \in \{0, 1\}$, outputs~$1$ if and only if one of the following
  conditions hold:
  \begin{align*}
    & (i_1, i_2 < L) \land (i_3, i_4, i_5 < R) \land
      \circuit(i_1, i_2, i_3, i_4, i_5, o_1, o_2, o_3, o_4, o_5) = 1\;,\\
    & (i_1 \geq L) \land (\text{$i_1$ is odd}) \land (o_1 = 1)\;,\\
    & (i_1 \geq L) \land (\text{$i_1$ is even}) \land (o_1 = 0)\;,\\
    & (i_2 \geq L) \land (\text{$i_2$ is odd}) \land (o_2 = 1)\;,\\
    & (i_2 \geq L) \land (\text{$i_2$ is even}) \land (o_2 = 0)\;.
  \end{align*}
  Now we show that $\circuit'$ succinctly describes $\decider$ on inputs~$n$,
  $x$, and~$y$ and time~$T$.
  To begin, we describe the decoupled 5SAT formula $\varphi_{\circuit'}$.
  The second and third constraints imply that for each $i_1 \geq L$,
  $\varphi_{\circuit'}$ contains the constraint $(a_{i_1})$ if $i_1$ is odd and
  $(\neg a_{i_1})$ if $i_1$ is even.
  Likewise, the fourth and fifth constraints imply that for each $i_2 \geq L$,
  $\varphi_{\circuit'}$ contains the constraint $(b_{i_2})$ if $i_2$ is odd and
  $(\neg a_{i_2})$ if $i_2$ is even.
  Thus, we can express the formula $\varphi_{\circuit'}$ as follows.
  \begin{equation}\label{eq:write-formula}
    \varphi_{\circuit'}(a, b, w_1, w_2, w_3)
    = \varphi_{\circuit}(a_1, b_1, w_{1, 1}, w_{2, 1} ,w_{3, 1})
  	\land (a_2 = (10)^{(L'-L)/2}) \land (b_2 = (10)^{(L'-L)/2}).
  \end{equation}
  Here, we write $a = (a_1, a_2)$ and $b = (b_1, b_2)$, where $a_1, b_1$ have
  length $L$ and $a_2, b_2$ have length $L' - L$, and for each $i \in \{1, 2,
  3\}$, we write $w_i = (w_{i, 1}, w_{i, 2})$, where $w_{i, 1}$ has length~$R$
  and $w_{i, 2}$ has length~$R'-R$.
  
  Now, suppose there exist $w_1, w_2, w_3$ such that $a, b, w_1, w_2, w_3$
  satisfy $\varphi_{\circuit'}$.
  Then because $\circuit$ succinctly represents~$\decider$, there exists
  $a_{\mathrm{prefix}}, b_{\mathrm{prefix}}$ of lengths $\ell_a, \ell_b \leq T$
  such that $\decider$ accepts $(n, x, y, a_{\mathrm{prefix}},
  b_{\mathrm{prefix}})$.
  In addition,
  \begin{equation*}
    a_1 = \mathrm{enc}_\Gamma(a_{\mathrm{prefix}}, \sqcup^{L/2-\ell_a}),
  \end{equation*}
  and likewise for~$b_1$.
  Equation~\eqref{eq:write-formula} then implies that
  \begin{align*}
    a = (a_1, a_2)
    &= (\mathrm{enc}_\Gamma(a_{\mathrm{prefix}}, \sqcup^{L/2-\ell_a}),
      (10)^{(L' - L)/2})\\
    &= (\mathrm{enc}_\Gamma(a_{\mathrm{prefix}}, \sqcup^{L/2-\ell_a}),
      \mathrm{enc}_\Gamma(\sqcup)^{(L' - L)/2})\\
    &= \mathrm{enc}_\Gamma(a_{\mathrm{prefix}}, \sqcup^{L'/2-\ell_a}),
  \end{align*}
  where the third step used the fact that $\mathrm{enc}_\Gamma(\sqcup) = 10$.
  As a similar statement holds for~$b$, this establishes that $\circuit'$
  succinctly describes $\decider$ on inputs~$n$, $x$, and~$y$ and time~$T$.
\end{proof}

By combining \Cref{prop:explicit-succinct-deciders} and
\Cref{prop:exciting-padding-prop} (choosing $\ell' = r' = r_0$ where $r_0$ is
from \Cref{prop:explicit-succinct-deciders}), we get the following:

\newcommand{\paddedsuccinctdecider}{\mathsf{PaddedSuccinctDecider}}

\begin{proposition}[Explicit padded succinct descriptions]
  \label{prop:explicit-padded-succinct-deciders}
  There is a Turing machine $\paddedsuccinctdecider$ with the following
  properties.
  Let $\decider$ be a decider, let $n$, $T$,~$\qlen$, and $\sigma$ be integers
  with $\qlen \leq T$ and $|\decider| \leq \sigma$, and let~$x$ and~$y$ be
  strings of length at most~$\qlen$.
  Then on input $(\decider, n, T, \qlen, \sigma, x, y)$,
  $\paddedsuccinctdecider$ outputs a circuit~$\circuit$ with five inputs of
  length~$m(T,\sigma)$ and five single-bit inputs which succinctly describes
  $\decider$ on inputs~$n$, $x$, and~$y$ and time~$T$.
  Moreover, the following hold.
  \begin{enumerate}
  \item $m(T,\sigma) = O(\log(T)+ \log(\sigma))$ and $2^m \geq 2T$.
  \item $\circuit$ has at most $s(n,T,\qlen,\sigma)$ gates, where $s$
    is such that $s(n,T,\qlen, \sigma) = \poly(\log(T), \log(n),
    \qlen,\sigma)$ and $5m(T, \sigma) + 5 + s(n, T, \qlen,\sigma)$ is
    a power of $2$.
  \item $\paddedsuccinctdecider$ runs in time $\poly(\log(T), \log(n),
    \qlen,\sigma)$, and the parameters $m(T,\sigma), s(n, T, \qlen, \sigma)$ can
    be computed from $n, T, \qlen,\sigma$ in time $\poly(\log(T), \log(n),
    \log(\qlen),\sigma)$.
  \end{enumerate}
\end{proposition}

\subsection{A PCP for normal form deciders}
\label{sec:pcp-cktval-new}

We give a probabilistically checkable proof (PCP) for the Bounded Halting
problem specialized to the case of normal form deciders.
Our PCP will use standard techniques from the algebraic, low-degree-code-based
PCP literature.
In particular, we slightly modify the PCP for Succinct-3SAT described
in~\cite[Section 11]{NW19} (which itself is based on the proof of the PCP
theorem in~\cite{harsha2004robust}) to apply it to the decoupled Succinct-5SAT
instances described in Section~\ref{sec:succinct-deciders}.
We follow their treatment closely.
As the PCPs constructed in this section are only an intermediate object towards
the normal form verifier introduced in the next section we do not include
standard definitions on PCPs, and refer to these references (in
particular~\cite{harsha2004robust}) for background.
We begin with some preliminaries.

\subsubsection{Preliminaries}

A key part of the PCP will be to design a function $f:\F^n \rightarrow \F$ which
is zero on the subcube $H_{\mathrm{subcube}} = \{0,1\}^n$.
Our next proposition shows that given such a function~$f$, there is a way of
writing it so that the fact that it is zero on $H_{\mathrm{subcube}}$ is
self-evidently true.
Doing so involves showing that $f$ can be written in a simple basis of
polynomials which are constructed to be zero on the subcube.
This fact is standard in the literature (see, for example,
\cite[Lemma~$4.11$]{ben2008short}), and we include its proof for
completeness.

\begin{proposition}[Polynomial basis of zero functions]\label{prop:zero-basis}
  Let $\F$ be a field.
  Define $\mathrm{zero}: \F \to \F$ as the univariate polynomial $x \mapsto
  x(1-x)$.
Define $H_{\mathrm{subcube}} = \{0,1\}^n$.
  Suppose $f:\F^n \rightarrow \F$ is an individual degree $d$ polynomial such
  that $f(x) = 0$ for all $x \in H_{\mathrm{subcube}}$.
  Then there exist polynomials $c_1, \ldots, c_n:\F^n \rightarrow \F$ such that
  for all $x \in \F^n$,
  \begin{equation*}
  f(x) = \sum_{i=1}^n c_i(x) \cdot \mathrm{zero}(x_i)\;.
  \end{equation*}
  In addition, for each $i \in \{1, \ldots, n\}$, $c_i$ has individual degree-$d$. 
\end{proposition}
  
  \begin{proof}
    To prove this, we first prove the following statement for each $k \in \{0,
    1, \ldots, n\}$: there exists an individual degree-$d$ polynomial $r_k:\F^{n}
    \rightarrow \F$ and polynomials $c_1, \ldots, c_k:\F^n \rightarrow \F$ such
    that
  \begin{equation}\label{eq:k-levels-of-division}
  f(x) = \sum_{i=1}^k c_i(x) \cdot \mathrm{zero}(x_i) + r_k(x)\;.
  \end{equation}
  In addition, for each $i \in \{1, \ldots, k\}$, $c_i$ has individual degree
  $d$ and the degree of~$x_i$ in $r_k$ is at most $1$.
  
  The proof is by induction on~$k$, the base case of $k = 0$ being trivial.
  Now, we perform the induction step.
  Assuming that Equation~\eqref{eq:k-levels-of-division} holds for $k$, we will
  show that it holds for~$k+1$ as well.
  Let~$r_k$ be the polynomial guaranteed by the inductive hypothesis.
  We now divide~$r_k$ by $\mathrm{zero}(x_{k+1})$ using polynomial
  division.
  This guarantees a polynomial $c_{k+1}$ and an individual degree-$d$
  polynomial~$r_{k+1}(x)$ such that
  \begin{equation*}
  r_k(x) = c_{k+1}(x) \cdot \mathrm{zero}(x_{k+1}) + r_{k+1}(x)\;.
  \end{equation*}
  In addition, $c_{k+1}$ still has individual degree $d$ (and in fact the degree
  of $x_{k+1}$ in $c_{k+1}$ is at most $d - 2$).
  Furthermore, for each $i \in \{1, \ldots, k+1\}$, $r_{k+1}$ has degree at most
  $1$ in~$x_i$.
  Plugging this into Equation~\eqref{eq:k-levels-of-division}, we see that
  \begin{equation}\label{eq:almost-there}
  f(x) = \sum_{i=1}^{k+1} c_i(x) \cdot \mathrm{zero}(x_i) + r_{k+1}(x)\;.  
  \end{equation}
  This completes the induction.
  
  Applying the $k = n$ case of Equation~\eqref{eq:k-levels-of-division}, we see
  that
  \begin{equation}\label{eq:case-n-of-induction}
  f(x) = \sum_{i=1}^n c_i(x) \cdot \mathrm{zero}(x_i) + r(x)\;,
  \end{equation}
  where the degree of each variables in $r$ is at most $1$.
  Since $f(x)$ vanishes on the subcube $H_{\mathrm{subcube}}$, it must also be
  that $r(x)$ vanishes on $H_{\mathrm{subcube}}$ as well.
  We claim that $r$ must therefore be the zero polynomial.
  
  We prove this by showing the following statement: for every polynomial
  $s(x_1,\ldots,x_n)$ with individual degree $1$ (i.e.
  a multilinear polynomial) that vanishes on all points $x \in \{0,1\}^n$, $s$
  must be the zero polynomial.
  We show this by induction on the number of variables.
  
  Consider the base case $n = 1$.
  Then $s$ is a univariate polynomial with degree at most $1$.
  However it is zero on two points, so therefore it must be the zero polynomial.
  Now for the inductive step: assuming the proposition holds for some $n \geq
  1$, we show that it holds for $n+1$ as well.
  For any multilinear polynomial$s$, we can write
	\[
		s(x_1,\ldots,x_{n+1}) = x_{n+1} s_1(x_1,\ldots,x_n) + s_2(x_1,\ldots,x_n)
	\]
	where both $s_1$ and $s_2$ are $n$-variate multilinear polynomials.
  Fix $x_{n+1} = 0$.
  Then $s(x_1,\ldots,x_n,0) = s_2(x_1,\ldots,x_n)$.
  Since $s$ vanishes on $\{0,1\}^{n+1}$, this implies that $s_2$ vanishes on
  $\{0,1\}^n$ as well.
  By the inductive hypothesis, $s_2$ is the identically zero polynomial.
  Now fix $x_{n+1} = 1$.
	Then $s(x_1,\ldots,x_n,1) = s_1(x_1,\ldots,x_n)$.
  Again, since $s$ vanishes on $\{0,1\}^{n+1}$, this implies that $s_1$ vanishes
  on $\{0,1\}^n$ as well.
  Again by the inductive hypothesis, $s_1$ is the identically zero polynomial.
  This shows that $s$ is the identically zero polynomial, completing the
  induction.
  


  Thus, $r$ is the zero polynomial.
  Applying this fact to Equation~\eqref{eq:case-n-of-induction},
  we arrive at the statement in the proposition.
  \end{proof}

\subsubsection{The PCP}
\label{sec:ar-pcp}

\paragraph{The problem.}
The input to the PCP verifier is a tuple $(\decider, n, T, \qlen, \sigma, \gamma,x,
y)$.
Here, $\decider$ is a decider, $n$, $T$, $\qlen$, $\sigma$ and $\gamma$ are integers
with $\qlen \leq T$ and $|\decider| \leq \sigma$, and~$x$ and~$y$ are a pair of
strings of length at most~$\qlen$ each.
The goal of the verifier is to check whether there exists two
strings~$a_{\mathrm{prefix}}$ and~$b_{\mathrm{prefix}}$ of length at most~$T$ 
such that $\decider$ halts on input
$(n,x,y,a_{\mathrm{prefix}},b_{\mathrm{prefix}})$ in time at most~$T$.
To do that the verifier makes random queries to a specially encoded \emph{PCP
  proof} $\Pi$, and decides whether to accept or reject based on the parts of
$\Pi$ that it reads. The proof $\Pi$ is encoded using a low-degree
code whose parameters are controlled by $\gamma$: the larger $\gamma$,
the smaller the soundness error in testing the code.
We first set the parameters used in the PCP construction.


\begin{definition}[Parameters for the PCP]\label{def:pcpparams}
  For all integers $n, T, \qlen, \sigma, \gamma
  \in \N$ such that $\qlen \leq T$ and $|\decider|\leq \sigma$ define the tuple
  $\pcpparams(n,T,\qlen,\sigma,\gamma) = (q,m,d,m',s)$ as follows.
  Let $m=m(T,\sigma)$, and $s=s(n,t,Q,\sigma)$ be as
  in~\Cref{prop:explicit-padded-succinct-deciders}.
  Let $a' > 1$ and $0 < b' < 1$ be the universal constants from
  \Cref{lem:ld-soundness}.
  Define the following integers.
  \begin{enumerate}
	\item Let $m' = 5m + 5 + s$ (recall that by
          \Cref{prop:explicit-padded-succinct-deciders}, $s$ is chosen
          so that $m'$ is a power of 2).
  \item Let $q = 2^k$ where $k$ is the smallest odd integer satisfying the following:
  \begin{enumerate}
  	\item $k \geq \Big ( (\gamma b' + 3a')/b' \Big ) \cdot \log s$.
	\item $(2 + 5k)m'/2^k < 1/2$.
	\item $km'/2^k \leq s^{-b' \gamma}$.
        \item $2^k$ is divisible by $m'$.
	\end{enumerate}
  \item Let $d = k$.

  \end{enumerate}
  Given $n, T, \qlen,\sigma,\gamma$ represented in binary, the parameter tuple
  $\pcpparams(n, T, \qlen, \sigma, \gamma)$ can be computed in time
  $\poly(\log(n), \log(T), \log(\qlen), \log (\gamma), \log(\sigma))$.
\end{definition}



Next, we define the format of a valid PCP proof, which for our construction
consists of evaluation tables of low-degree polynomials.

\begin{definition}\label{def:pcp-proof}
  Given $n, T, \qlen, \sigma, \gamma \in \N$ and $(q,m,d,m',s)=
  \pcpparams(n, T, \qlen, \sigma,\gamma)$, a \emph{low-degree PCP proof} is a tuple
  $\Pi$ of evaluation tables of polynomials $g_1, \dots, g_5: \F_q^{m} \to \F_q$
  and $c_0, \dots, c_{m'}: \F_q^{m'} \to \F_q$ with all polynomials having
  individual degree at most $d$.
  We divide the $m'$ input variables of $c_0, \dots, c_{m'}$ into blocks as
  follows:
  \[
    \F_q^{m'} \ni z = (\underbrace{x_1}_{\F_q^m}, \dots,
    \underbrace{x_5}_{\F_q^m}, \underbrace{o}_{\F_q^5},
    \underbrace{w}_{\F_q^s})\;.
  \]
\end{definition}

\begin{definition}\label{def:pcp-eval}
  Given a low-degree PCP proof $\Pi$ and a point $z = (x_1, \dots, x_5, o, w)
  \in \F_q^{m'}$, where $x_1, \dots, x_5 \in \F_q^{m}$, $o \in \F_q^{5}$, and $w
  \in \F_q^{s}$, the \emph{evaluation} of $\Pi$ at $z$ is given by
  \[
    \ev_z(\Pi) = (\alpha_1, \dots, \alpha_5, \beta_0, \dots, \beta_{m'}) \in
    \F_{q}^{6 + m'}\;,
  \]
  where $\alpha_i = g_i(x_i)$ and $\beta_j = c_j(z)$.
\end{definition}

\begin{theorem}\label{thm:pcp-decider}
  There exists a Turing machine $\pcpverifier$ with the following properties.
  \begin{enumerate}
  \item \textbf{(Input format)} The input to $\pcpverifier$ consists of two parts:
    a ``decider specification'' and a ``PCP view.''
    \begin{enumerate}
    \item \textbf{(Decider specification)} Let $\decider$ be a decider, $n$,
      $T$, $\qlen$, $\sigma$ and $\gamma$ be integers with $\qlen \leq T$, and
      $\sigma = |\decider|$, and let~$x$ and~$y$ be strings of length at
      most~$\qlen$.
      Let $(q,m,d,m',s) = \pcpparams(n, T, \qlen,\sigma, \gamma)$ be as in
      \Cref{def:pcpparams}.
      Then the corresponding decider specification is the tuple
      \[
        (\decider, n, T, \qlen, \sigma, \gamma, x, y).
      \]
    \item \textbf{(PCP view)} Let $z \in \F_{q}^{m'}$ and let $\pcpeval \in
      \F_q^{6 + m'}$.
      Then the PCP view is the pair $(z, \pcpeval)$.
    \end{enumerate}
    The Turing machine $\pcpverifier$ returns either $1$ (accept) or $0$
    (reject).
  \item
    \textbf{(Complexity):} 
$\pcpverifier$ runs in time at most~$\poly(\log(T), \log(n), \qlen, \sigma, \gamma)$.
        
  \end{enumerate}
  For the remaining items, fix a decider specification $(\decider, n, T, \qlen,
  \sigma, \gamma, x, y)$; we think of $\pcpverifier$ as a function of the PCP view
  input only.

  \begin{enumerate}
  \item[3.]
    \textbf{(Completeness):} Suppose $a_{\mathrm{prefix}}, b_{\mathrm{prefix}}
    \in \{0, 1\}^*$ are two strings of length $\ell_a, \ell_b \leq T$,
    respectively, such that $\decider$ halts and accepts in time~$T$ on input $(n, x, y,
    a_{\mathrm{prefix}}, b_{\mathrm{prefix}})$.
    Let $M = 2^m$ and write
    \begin{equation*}
      a = \mathrm{enc}_\Gamma(a_{\mathrm{prefix}}\,,\, \sqcup^{M/2 - \ell_a})
      \quad
      \text{and}
      \quad
      b = \mathrm{enc}_\Gamma(b_{\mathrm{prefix}}\,,\, \sqcup^{M/2 - \ell_b})\;,
    \end{equation*}
    where recall that by
    \cref{prop:explicit-padded-succinct-deciders}, $m$ is chosen so that $2T\leq M$, and $\sqcup$ is a special symbol that
    is encoded using two bits; thus $a$ and $b$ are each an $M$-bit string.
    Then there exists a low-degree PCP proof (Definition~\ref{def:pcp-proof})
    $\Pi = (g_1, \ldots, g_5, c_0, \ldots, c_{m'})$ with $g_1 = g_a$ and $g_2 =
    g_b$, the low-degree encodings of $a$ and $b$, respectively (see
    \Cref{sec:ld-encoding}), which causes $\pcpverifier$ to accept with
    probability~$1$ over a uniformly random $z \in \F_q^{m'}$:
    \[ \Pr_{z \in \F_q^{m'}}\big(\pcpverifier(z, \ev_{z}(\Pi)) = 1\big) = 1\;. \]
  \item[4.]
    \textbf{(Soundness):} Let $\Pi = (g_1, \ldots, g_5, c_0, \ldots, c_{m'})$ be a
    low-degree PCP proof such that $\pcpverifier$ at a uniformly random $z$
    accepts with probability larger than $p_{\soundness} =\frac{1}{2}$:
    \[ \Pr_{z \in \F_q^{m'}}\big( \pcpverifier(z, \ev_{z}(\Pi)) = 1\big) >
      p_{\soundness}\;. \] Then there exist strings $a, b \in \{0, 1\}^{M}$ with the
    following properties.
    \begin{enumerate}
    \item There exist strings $a_{\mathrm{prefix}}, b_{\mathrm{prefix}} \in \{0,
      1\}^*$ of length $\ell_a, \ell_b$, respectively, such that
      \begin{equation*}
        a = \mathrm{enc}_\Gamma(a_{\mathrm{prefix}}\,,\, \sqcup^{M/2 - \ell_a})
        \quad
        \text{and}
        \quad
        b = \mathrm{enc}_\Gamma(b_{\mathrm{prefix}}\,,\, \sqcup^{M/2 - \ell_b})\;.
      \end{equation*}
    \item $\decider$ halts in time~$T$ and accepts on input $(n, x, y, a_{\mathrm{prefix}},
      b_{\mathrm{prefix}})$.
    \item $a = \coded(g_1)$ and $b = \coded(g_2)$, where $\coded(\cdot)$ is the
      Boolean decoding map of the low-degree encoding defined in
      \Cref{sec:ld-encoding}.
\end{enumerate}

  \end{enumerate}
\end{theorem}
It is important to note that the soundness in~\Cref{thm:pcp-decider} is
\emph{only} against low-degree proofs, which is why we are able to
obtain a soundness error independent of $\gamma$. The eventual
answer-reduced verifier (\Cref{thm:ar}) will combine this PCP with the low-degree
test, and the latter will introduce a dependence on $\gamma$ in the
soundness error.

\begin{proof}[Proof of \cref{thm:pcp-decider}]


  We present the description of the Turing machine $\pcpverifier$ in
  \Cref{fig:pcpverifier}.
  
\begin{figure}[!htb]
  \begin{gamespec}
  The Turing machine $\pcpverifier$ takes the following as input:
	\begin{itemize}
		\item Decider specification $(\decider, n, T, \qlen,
      \sigma, \gamma, x, y)$
      	\item PCP view $(z,\Xi)$
	\end{itemize}
	The Turing machine performs the following steps sequentially:
  \begin{enumerate}
  \item \label{enu:compute-circuit} Compute $\circuit =
    \paddedsuccinctdecider(\decider, n, T, \qlen, \sigma, x, y)$.
    The circuit $\circuit$ has five $m$-bit inputs and five single-bit inputs,
    and contains at most~$s$ AND and OR gates.
    


  \item \label{enu:compute-tseitin-formula} Compute the Tseitin formula
    $\formula$ (\cref{def:tseitin}) corresponding to $\circuit$, which is a boolean formula on $m' =
    5m + 5 + s$ variables.
    Let $\formula_{\mathrm{arith}}: \F_q^{m'} \to \F_q$ denote the
    arithmetization of $\formula$ (\cref{def:formula-arithmetization}).
  \item \label{enu:parse-input} Parse $z = (x,o,w) \in \F_q^{m'}$ where
    $x_1,\ldots,x_5 \in \F_q^m$, $o \in \F_q^5$, and $w \in \F_q^s$. (Here $x$ is a variable which should not be confused with an input $x$ to $\decider$.)
    Parse $\Xi = (\alpha_1,\ldots,\alpha_5,\beta_0,\ldots,\beta_{m'}) \in
    \F_q^{6 + m'}$.
  
  \item (\textbf{Formula test}) \label{enu:formula-test} Reject if $\beta_0 \neq
    \formula_{\mathrm{arith}}(x,o,w) \cdot (\alpha_1 - o_1) \cdots (\alpha_5 -
    o_5)$.
    Otherwise, continue.
  \item (\textbf{Zero on subcube test}) \label{enu:zero-test} Reject if $\beta_0
    \neq \sum_{i=1}^{m'} \beta_i \cdot \mathrm{zero}(z_i)$.
    Otherwise, accept.
  \end{enumerate}
  \end{gamespec}
\caption{The decision procedure $\pcpverifier$.}
\label{fig:pcpverifier}
\end{figure}

Before establishing the Complexity, Completeness, and Soundness properties of
$\pcpverifier$, we point out several important items.
Recall that $M = 2^m$.
\begin{enumerate}
\item We recall what it means for the circuit $\circuit$ computed in
  Step~\ref{enu:compute-circuit} of \Cref{fig:pcpverifier} to succinctly
  describe the 5SAT formula $\varphi_\circuit$.
  Let the inputs of the formula $\varphi_\circuit$ be strings $a,b,u_3,u_4,u_5
  \in \{0,1\}^M$.
  Then for all $x_1,\ldots,x_5 \in \{0,1\}^m$ and $o \in \{0,1\}^5$, we have
  that $\circuit(x_1,\ldots,x_5,o) = 1$ if and only if $a_{x_1}^{o_1} \lor
  b_{x_2}^{o_2} \lor u_{3, x_3}^{o_3} \lor u_{4, x_4}^{o_4} \lor u_{5,
    x_5}^{o_5}$ is a clause in~$\varphi_\circuit$, where the coordinates of
  $a,b,u_3,u_4,u_5$ are indexed by strings of length $m$.

  Furthermore, the formula $\varphi_\circuit$ is related to the decider
  $\decider$ in the following way.
  For all $a, b\in \{0, 1\}^M$, there exist $u_3, u_4, u_5 \in \{0, 1\}^M$ such
  that $a, b, u_3, u_4, u_5$ satisfy $\varphi_{\circuit}$ if and only if there
  exist $a_{\mathrm{prefix}}, b_{\mathrm{prefix}} \in \{0, 1\}^*$ of
  lengths~$\ell_a, \ell_b \leq T$, respectively, such that
  \begin{equation*}
    a = \mathrm{enc}_\Gamma(a_{\mathrm{prefix}}, \sqcup^{M/2 - \ell_a})
    \quad \text{and} \quad
    b = \mathrm{enc}_\Gamma(b_{\mathrm{prefix}}, \sqcup^{M/2 - \ell_b})
  \end{equation*}
  and $\decider$ accepts $(n, x, y, a_{\mathrm{prefix}}, b_{\mathrm{prefix}})$
  in time~$T$.
  
\item We recall from \Cref{def:tseitin} that the boolean formula $\formula$ is related to $\circuit$ in the following
  way.
  For all $x_1, \ldots,x_5 \in \{0, 1\}^m$ and $o \in \{0, 1\}^5$,
  $\circuit(x_1,\ldots,x_5, o) = 1$ if and only if there exists a $w \in \{0,
  1\}^s$ such that $\formula(x_1,\ldots,x_5, o, w) = 1$.
  The size of the formula $\formula$ is linear in the size of the circuit
  $\circuit$, which is $s$.
  
\item We recall from \Cref{def:formula-arithmetization} that the arithmetization $\formula_{\mathrm{arith}}:=
  \mathrm{arith}_q(\formula)$ is a function
  $\formula_{\mathrm{arith}}:\F_q^{m'}\rightarrow \F_q$ such that
  \begin{equation}
    \label{eq:farith}
    \forall (x,o,w) \in \{0, 1\}^{m'}\;,
    \quad \formula_{\mathrm{arith}}(x,o,w) = \formula(x,o,w)\;.
  \end{equation}
  By \cref{prop:tseitin-arith-degree}, $\formula_{\mathrm{arith}}$ is an
  individual degree $2$ polynomial.
\end{enumerate}

\paragraph{Complexity.}
We bound the complexity of the Turing machine $\pcpverifier$.
From \Cref{prop:explicit-padded-succinct-deciders},
Step~\ref{enu:compute-circuit} takes time $\poly(\log(T),\log(n),Q,\sigma)$.
Computing the Tseitin formula in Step~\ref{enu:compute-tseitin-formula} takes
time that is polynomial in the size of the circuit $\circuit$, which is
$\poly(s) = \poly(\log(T),\log(n),Q,\sigma)$.
The Formula Test, Step~\ref{enu:formula-test}, requires computing the
arithmetization $\formula_{\mathrm{arith}}$ at a point $z \in \F_q^{m'}$, which
takes time $\poly(s,\log q)$.
The Zero on Subcube Test, Step~\ref{enu:zero-test}, takes time $\poly(m',\log
q)$ to compute a sum of $m'$ products of $\F_q$ elements (evaluating
$\mathrm{zero}(\cdot)$ takes time $\poly(\log q)$).
Thus the complexity of $\pcpverifier$ is bounded by
\[
  \poly(\log(T),\log(n),Q,\sigma,\gamma)
\]
using our choice of $\pcpparams$ from \Cref{def:pcpparams}. (The
dependence on $\gamma$ comes form the setting of $q$.) 
  
  
\paragraph{Completeness.}
We now establish the Completeness property of $\pcpverifier$ by concocting a PCP
proof $\Pi$ that is accepted with probability $1$.
  
Let $a,b \in \{0,1\}^M$ be as specified in Item 3 of \Cref{thm:pcp-decider}.
Then by definition of the formula $\varphi_\circuit$ and the assumption that 
 $\decider$ halts and accepts in time~$T$ on input $(n, x, y,
    a_{\mathrm{prefix}}, b_{\mathrm{prefix}})$, there exist strings
$u_3,u_4,u_5 \in \{0,1\}^M$ such that $(a,b,u_3,u_4,u_5)$ satisfies
$\varphi_\circuit$.
  
Let $g_1,g_2$ denote the low-degree encodings of $a$ and $b$, respectively (see
\Cref{sec:ld-encoding} for definition of low-degree encoding), and let
$g_3, g_4, g_5$ denote the low-degree encodings of $u_3, u_4$, and
$u_5$, respectively.
In particular, $g_1,\ldots,g_5$ are $m$-variate multilinear polynomials, and for
all $x_1,\ldots,x_5 \in \{0,1\}^m$,
\begin{equation*}
  g_1(x_1) = a_{x_1}\;, \quad
  g_2(x_2) = b_{x_2}\;, \quad
  g_3(x_3) = u_{3,x_3}\;, \quad
  g_4(x_4) = u_{4,x_4}\;, \quad
  g_5(x_5) = u_{5,x_5}\;,
\end{equation*}
where we index the coordinates of $a,b,u_3,u_4,u_5$ by $m$-bit strings in the
natural way.

Next, define the polynomial $c_0:\F_q^{m'} \rightarrow \F_q$ as
\begin{equation*}
  c_0(x, o, w) = \formula_{\mathrm{arith}} (x, o, w)
  \cdot (g_1(x_1) - o_1) \cdots (g_5(x_5)- o_5)\;.
\end{equation*}
Note that $c_0$ has individual degree $3$, since as recalled in item 3 above
$\formula_{\mathrm{arith}} (x, o, w)$ has individual degree $2$ and each of the
$g_i$'s are multilinear).
  
We now show that $c_0(x,o,w) = 0$ for all $(x,o,w) \in \{0,1\}^{m'}$.
By construction, the arithmetization $\formula_{\mathrm{arith}} (x, o, w)$ is
either $0$ or $1$.
If it is $0$, then we are done.
If it is $1$, then this means that $\formula(x,o,w) = 1$, which by definition
means that $\circuit(x,o) = 1$, which means that $a_{x_1}^{o_1} \lor
b_{x_2}^{o_2} \lor u_{3, x_3}^{o_3} \lor u_{4, x_4}^{o_4} \lor u_{5, x_5}^{o_5}$
is a clause in the formula $\varphi_\circuit$.
But since $(a,b,u_3,u_4,u_5)$ satisfies $\varphi_\circuit$, this means that at
least one of $a_{x_1} = o_1$, $b_{x_2} = o_2$, $u_{3,x_3} = o_3$, $u_{4,x_4} =
o_4$, or $u_{5,x_5} = o_5$.
But this means that the product $(g_1(x_1) - o_1) \cdots (g_5(x_5)- o_5) = 0$,
by definition of the $g_i$'s.
Thus $c_0(x,o,w) = 0$.
    
Thus, by \Cref{prop:zero-basis}, there exist polynomials $c_1,\ldots,c_{m'}:
\F_q^{m'} \to \F_q$ of individual degree at most $3$ that certify that $c_0$
vanishes on $\{0,1\}^{m'}$.
In other words, for all $z \in \F_q^m$ (not just over the boolean cube), we have
\begin{equation}
	\label{eq:restricted-to-subcube}
  c_0(z) = \sum_{i=1}^{m'} c_i(z) \cdot \mathrm{zero}(z_i)
\end{equation}
where $\mathrm{zero}(x) = x(1-x)$.
    
Let the PCP proof $\Pi$ be the collection of polynomials
$(g_1,\ldots,g_5,c_0,c_1,\ldots,c_{m'})$, where the $g_i$'s act on disjoint
variables and the $c_i$'s act on all of them.
Let $z \in \F_q^{m'}$ and let $\Xi =
(\alpha_1,\ldots,\alpha_5,\beta_0,\beta_1,\ldots,\beta_{m'}) = \eval_z(\Pi)$.
The PCP view $(z,\Xi)$ passes the Formula Test always, because $\beta_0 =
c_0(z)$ and $\alpha_i = g_i(z)$ for $i \in \{1,2,\ldots,5\}$.
The PCP view also passes the Zero on Subcube Test, because of
\Cref{eq:restricted-to-subcube}.
    
Thus the PCP view is accepted by the Turing machine $\pcpverifier$ with
probability $1$.
This establishes the Completeness property.
    
\paragraph{Soundness.}
Fix a low-degree PCP proof $\Pi = (g_1,\ldots,g_5,c_0,\ldots,c_{m'})$ such that
the PCP view $(z,\Xi)$ for $\Xi = \eval_z(\Pi)$ causes $\pcpverifier$ to accept
with probability greater than $p_\soundness$.
This, in particular, implies that the Formula Test is satisfied with high
probability.
With probability greater than $p_\soundness$ over the choice of $z \sim
\F_q^{m'}$, we have
\[
  c_0(z) = \formula_{\mathrm{arith}} (z) \cdot (g_1(x_1) - o_1) \cdots (g_5(x_5)- o_5)
\]
Note that $c_0(z)$, by definition of low-degree PCP proof, has individual degree
at most $d$.
The right-hand side has individual degree at most $2 + 5d$, because
$\formula_{\mathrm{arith}}$ has individual degree $2$ and the $g_i$'s have
individual degree at most $d$.
Thus both sides have total degree at most $(2 + 5d)m'$.
Suppose $c_0(z)$ was not identical to $\formula_{\mathrm{arith}} (z) \cdot
(g_1(x_1) - o_1) \cdots (g_5(x_5)- o_5)$.
Then the Schwartz-Zippel lemma implies that the probability they agree on a
randomly chosen $z$ is at most $(2 + 5d)m'/q$.
By our choice of $q$, $m'$, and $d$ in \Cref{def:pcpparams}, this is less than
$p_\soundness$, which is a contradiction.
Thus $c_0(z)$ is equal to $\formula_{\mathrm{arith}} (z) \cdot (g_1(x_1) - o_1)
\cdots (g_5(x_5)- o_5)$ for all $z\in \F_q^{m'}$.
    
Next, we consider the Zero on Subcube Test.
With probability greater than $p_\soundness$ over the choice of $z \sim
\F_q^{m'}$, we have
\[
  c_0(z) = \sum_{i=1}^{m'} c_i(z) \cdot \mathrm{zero}(z_i).
\]
Again, the polynomials on both sides of the equation have individual degree at
most $d+2$, and therefore total degree at most $(2 + d)m'$.
By the Schwartz-Zippel lemma, the probability that both sides would agree on the
evaluation of a random $z$, if they weren't equal polynomials, would be at most
$(2 + d)m'/q$, which is also less than $p_\soundness$.
Thus $c_0(z) = \sum_{i=1}^{m'} c_i(z) \cdot \mathrm{zero}(z_i)$ for all $z \in
\F_q^{m'}$.
This in particular implies that $c_0(z)$ vanishes on the subcube $\{0,1\}^{m'}$. 
    
We can now decode assignments $a,b,u_3,u_4,u_5 \in \{0,1\}^M$ that satisfy the
formula $\varphi_\circuit$.
Let $a = \coded(g_1), b = \coded(g_2), u_3 = \coded(g_3), u_4 = \coded(g_4), u_5
= \coded(g_5)$ where $\coded(\cdot)$ is the decoding map defined in
\Cref{sec:ld-encoding}.
The fact that $c_0(z)$ vanishes on the subcube $\{0,1\}^{m'}$ implies that all
clauses of $\varphi_\circuit$ are satisfied by $a,b,u_3,u_4,u_5$.
By construction of $\varphi_\circuit$, this implies that there exists
$a_{\mathrm{prefix}}, b_{\mathrm{prefix}}$ such that
$\decider(n,x,y,a_{\mathrm{prefix}},b_{\mathrm{prefix}})$ accepts in time $T$.
This completes the proof of the Soundness property, and the proof of the
Theorem.
\end{proof}

\subsection{A normal form verifier for the PCP}
\label{sec:ld-compiler}

\def\tvans{\hat{\verifier}^\ar}
\def\tsans{\hat{\sampler}^\ar}
\def\tdans{\hat{\decider}^\ar}
\def\vans{\verifier^\ar}
\def\sans{\sampler^\ar}
\def\dans{\decider^\ar}

In this section we show how to convert the PCP from \cref{sec:pcp-cktval-new}
into a normal form verifier.
This results in an ``answer reduction'' scheme: a way to map a verifier
$\verifier$ into a new (typed) verifier $\tvans$ with a smaller answer size.
	
Let $\verifier = (\sampler, \decider)$ be a normal form verifier and
$(\lambda,\mu,\gamma)$ a tuple of integers.
In the rest of this section we define the (typed) answer-reduced verifier $\tvans =
(\tsans, \tdans)$ associated with $\verifier$ and parameters $(\lambda, \mu, \gamma)$.
Completeness, complexity and soundness of the construction are shown in the
following sections.

\subsubsection{Parameters and notation}\label{sec:ar-params}
We establish some notation and parameters that are used throughout \Cref{sec:ld-compiler}. 
First, define the functions
\begin{equation}
  \label{eq:ar-params-1}
  T(n)=(2^{\lambda n})^\mu \quad\text{and} \quad \qlen(n) = (\lambda n)^{\mu}\;.
\end{equation} 
In the main theorem of this section, \Cref{thm:ar}, we assume that the input
verifier $\verifier$ satisfies
\begin{equation}
  \label{eq:ar-time-assumption}
  \TIME_\decider(n) \leq T(n)\quad
  \text{and}\quad \TIME_\sampler(n) \leq Q(n) \;.
\end{equation}
In what follows we write $T$ and $Q$ as free parameters (though they are
implicitly functions of the index $n$).

Next, for all integers $n \in \N$ define the PCP parameters $(q,m,d,m',s) =
\pcpparams(n, T, \qlen, \sigma, \gamma)$ where $\sigma = |\decider|$.
Note that the parameters $(q,m,d,m',s)$ are all implicitly functions of $n$.





Recall from Definition~\ref{def:sampler-sample} the notation $\mu_{\sampler}$
denoting the distribution over pairs of questions $(x_\alice,x_\bob)$ generated by
$\sampler$.
We use $\mu_{\sampler, \alice}$ to indicate the marginal distribution of
$\mu_\sampler$ on the first question $x_\alice$, and $\supp(\mu_{\sampler})$ to
indicate the set of question pairs that have nonzero probability under
$\mu_\sampler$.

\begin{remark}\label{rk:ab-01}
  Throughout \Cref{sec:ld-compiler}, for convenience we often identify the label
  $\alice$ with $1$ and $\bob$ with $2$.
  For example, $g_\alice$ is another label for the polynomial $g_1$.
\end{remark}

\subsubsection{The answer-reduced verifier}
\label{sec:ar-verifier}

Let $\verifier = (\sampler, \decider)$ be a normal form verifier and
$(\lambda,\mu,\gamma)$ integers.
All other required parameters and notation are introduced in
Section~\ref{sec:ar-params}.

\newcommand{\tspcp}{\hat{\sampler}^\pcp}

\paragraph{Sampler.}
In this section we define the (typed) sampler $\tsans$ for the (typed) answer reduced
verifier $\tvans$. 

We first give an intuitive description of the sampler, and then proceed to give
a formal definition.
The sampler $\tsans$ is a \emph{product} of the oracularized sampler $\tsora$
corresponding to $\sampler$, and a typed PCP sampler $\tspcp$.
The corresponding question distribution $\mu_{\tsans}$ is then a product
distribution $\mu_{\tsora} \times \mu_{\tspcp}$.

The PCP distribution $\mu_{\tspcp}$ corresponds to six copies of the classical
low-degree test question distribution (see \Cref{sec:ld-game}).
Recall that the classical low-degree test is a nonlocal game that checks whether
the players' answers are consistent with low-degree polynomials.
The answer-reduced verifier will use the classical low-degree test to check that
the players respond with evaluations of low-degree polynomials
$g_1,\ldots,g_5,c_0,\ldots,c_{m'}$, and then process the evaluations using the
PCP verifier $\pcpverifier$ from \Cref{thm:pcp-decider}.
Five of the six copies are meant to certify the ``low-degreeness'' of
$g_1,\ldots,g_5$, and the sixth copy is used to certify the ``low-degreeness''
of the entire bundle of polynomials $g_1,\ldots,g_5,c_0,\ldots,c_{m'}$.
The reason that the $g_i$'s are checked individually is to ensure that they only
depend on certain blocks of input variables, because ultimately the soundness of
the answer-reduced verifier reduces to the soundness of the PCP verifier
$\pcpverifier$, and \Cref{thm:pcp-decider} assumes that the polynomials
$g_1,\ldots,g_5$ of the PCP proof only depend on certain subsets of variables.

We now formally define the typed PCP sampler $\tspcp$. The type set is
\begin{equation*}
  \type^\pcp = \{ \Point_1, \ldots,\Point_6\} \cup \{\ALine_1,
  \dots, \ALine_6\} \cup \{\DLine_1,
  \dots, \DLine_6\}\;,
\end{equation*}
and the type graph $G^\pcp=(\type^\pcp,E^\pcp)$ uses the complete edge set
$E^\pcp = \type^\pcp \times \type^\pcp$.
The ambient vector space for the sampler is
\begin{equation}
  \label{eq:V-pcp}
  V^\pcp = \Bigl( \bigoplus_{i=1}^{5} V_{i, \xpt} \oplus V_{i, \coord} \oplus
  V_{i, \dir{}}\Bigr) \oplus V_{\aux, \xpt} \oplus V_{\aux, \coord} \oplus
  V_{\aux, \dir{}}\;,
\end{equation}
where the spaces $V_{i,\xpt}, V_{i, \dir{}}$ are each isomorphic to $\F_q^m$,
the space $V_{i, \coord}$ is isomorphic to $\F_{q}$, the spaces $V_{\aux, \xpt},
V_{\aux, \dir{}}$ are each isomorphic to $\F_q^{5+s}$, and the space $V_{\aux,
  \coord}$ is isomorphic to $\F_q$.
In addition, define the following direct sums:
\begin{align*}
  V_{6, \xpt} & = \Big( \bigoplus_{i=1}^{5} V_{i, \xpt} \Big) \oplus V_{\aux, \xpt}\;,\\
  V_{6, \coord} & = \Big( \bigoplus_{i=1}^{5} V_{i, \coord} \Big) \oplus
                   V_{\aux, \coord}\;,\\
  V_{6, \dir{}} & = \Big( \bigoplus_{i=1}^{5} V_{i, \dir{}} \Big) \oplus
                   V_{\aux, \dir{}}\;.
\end{align*}
For all $i \in \{1,2,\ldots,6\}$, we call the space $V_{i,\xpt}$ the
\emph{$i$-th point register}, the space $V_{i,\coord}$ the \emph{$i$-th
  coordinate register}, and the space $V_{i,\dir{}}$ the \emph{$i$-th direction
  register} (compare this with the subspaces defined in \Cref{sec:ld-game}).

For every type $\tvar \in \type^\pcp$, we define the following CL functions on
$V^\pcp$:
\begin{itemize}
\item For the types $\tvar = \Point_i$ for $i \in \{1, \dots, 6\}$, the
  corresponding $1$-level CL function $L_{\Point_i}$ is identical to the CL
  function $L_\Point$ defined in \Cref{eq:cl-ptf}, but acts on the
  subspace $V_{i,\xpt} \oplus V_{i, \coord} \oplus V_{i,
    \dir{}}$ (and zeroes out all the other registers).

\item For the types $\tvar = \ALine_i$ for $i \in \{ 1, \dots, 6\}$, the
  corresponding $2$-level CL function $L_{\ALine_i}$ is identical to the CL
  function $L_\ALine$ defined in \Cref{eq:cl-alnf}, but acts on the subspaces
  $V_{i,\xpt} \oplus V_{i,\coord} \oplus V_{i,\dir{}}$.
  
\item For the types $\tvar = \DLine_i$ for $i \in \{ 1, \dots, 6\}$, the
  corresponding $3$-level CL function $L_{\DLine_i}$ is identical to the CL
  function $L_\DLine$ defined in \Cref{eq:cl-dlnf}, but acts on the subspaces
  $V_{i,\xpt} \oplus V_{i,\coord} \oplus V_{i,\dir{}}$.
\end{itemize}

Observe that for $i \in \{1,\ldots,5\}$, the CL distributions
$\mu_{L_{\ALine_i},L_{\Point_i}}$ and $\mu_{L_{\DLine_i},L_{\Point_i}}$ correspond to
the question distributions of the classical low-degree test parameterized by $q$
and $m$ (see \Cref{sec:ld-game} for the definition of the low-degree test
distributions, which only depend on the first two arguments of the $4$-tuple
$\ldparams$), and the CL distributions $\mu_{L_{\ALine_6},L_{\Point_6}}$ and
$\mu_{L_{\DLine_6},L_{\Point_6}}$ correspond to the classical low-degree test
parameterized by $q$ and $m'$.

We finally give the formal definition of the sampler $\tsans$.
The type set of $\tsans$ is $\type^\ar = \type^\ora \times \type^\pcp$ and the
type graph is $G^\ar = G^\ora \times G^\pcp$ with edge set
\begin{equation*}
  E^\ar = \bigl\{\{(\lvar,\rvar),(\lvar',\rvar')\}:\,
  \{\lvar, \lvar'\} \in E^\ora \text{ and } \{\rvar, \rvar'\} \in
  E^\pcp \bigr\}\;.
\end{equation*}
Let $V^\ora$ denote the ambient space of the oracularized sampler $\tsora$.
The ambient space of $\tsans$ is then the direct sum $V^\ora \oplus V^\pcp$.
For all type pairs $(\tvar_\ora,\tvar_\pcp) \in \type^\ar$, the corresponding CL
function $L_{\tvar_\ora,\tvar_\pcp}$ is simply the direct sum of the CL functions $L_{\tvar_\ora}$ (coming from $\tsora$) and
$L_{\tvar_\pcp}$ (coming from $\tspcp$).
Thus one can see that the distribution corresponding to $\tsans$ is the product
distribution $\mu_{\tsora} \times \mu_{\tspcp}$.

Note that $\tsans$ is an $\max\{ \ell, 3\}$-level sampler, where $\ell$ is the
number of levels of the sampler $\sampler$.







\paragraph{Decider.}
The decider $\tdans$ is described in \cref{fig:decider-pcp}.

\begin{figure}[!htb]
  \begin{gamespec}
  {
  \small
    \begin{table}[H]
      \centering
      \small
      \begin{tabularx}{\textwidth}{ l l X }
        \toprule
        Type & Question Format & Answer Format \\
        \midrule
        $\Point_i$ for $i=1, \ldots, 5$ & $y_i \in \F_q^m$ &
        $\alpha_i \in \F_q$ \\
        $\ALine_i$ for $i=1, \ldots, 5$ & $v_i \in \F_q^m \times \F_q$ &
        $h_i: \F_q \to \F_q$ \\
        $\DLine_i$ for $i=1, \ldots, 5$ & $v_i \in \F_q^m \times \F_q \times \F_q^m$ &
        $h_i: \F_q \to \F_q$ \\        
        $\Point_6$ & $z = (y, o, w) \in \F_q^{m'}$ & $(\alpha'_1, \ldots,
        \alpha'_5, \beta_0, \ldots, \beta_{m'}) \in \F_q^{m'+6}$\\
        $\ALine_6$ & $v \in \F_q^{m'} \times \F_q$ & $(h'_1, \dots,
        h'_5, f_0, \dots, f_{m'}): \F_q \to \F_q^{m' + 6}$ \\
        $\DLine_6$ & $v \in \F_q^{m'} \times \F_q \times \F_q^{m'}$ & $(h'_1, \dots,
        h'_5, f_0, \dots, f_{m'}): \F_q \to \F_q^{m' + 6}$ \\        
        \bottomrule
      \end{tabularx}
      \caption{Question and answer formats for types in $\type^\pcp$.}\label{table:tpcp}
    \end{table}

    On input $(n, \tvar_\alice, x_\alice, \tvar_\bob, x_\bob, a_\alice,
    a_\bob)$, the decider $\tdans$ parses $\tvar_\alice$, $\tvar_\bob$ as
    $(\tvar_{Q,\alice}, \tvar_{\Pi, \alice})$, and $(\tvar_{Q,\bob},
    \tvar_{\Pi,\bob})$ respectively in $ \type^\ora \times \type^\pcp$, parses
    $x_\alice$ and $x_\bob$ as $(x_{Q,\alice}, x_{\Pi,\alice})$ and
    $(x_{Q,\bob}, x_{\Pi,\bob})$ respectively. The answer format depends only on $\tvar_{\Pi, \alice}$ and $\tvar_{\Pi, \bob}$ respectively and is as indicated in Table~\ref{table:tpcp}.  
    The decider performs the following steps sequentially, for all
    $w\in\{\alice,\bob\}$:
    \begin{enumerate}
    \item (\textbf{Global consistency check}): If $\tvar_{\alice} =
      \tvar_{\bob}$, reject if $a_\alice\neq a_\bob$.
      \label{enu:ar-global-consistency}

    \item (\textbf{Input consistency check}): If $\tvar_{Q,w} = \oracle$ and
      $\tvar_{Q, \ol{w}} = v \in \{\alice,\bob\}$, and if $(\tvar_{\Pi,w},
      \tvar_{\Pi, \ol{w}}) = (\Point_6, \Point_{v})$, reject if $\alpha_{v} \neq
      \alpha'_{v}$ (where $\alice \leftrightarrow 1$ and $\bob \leftrightarrow
      2$, as per Remark~\ref{rk:ab-01}).
      \label{enu:ar-input-consistency}

    \item (\textbf{Input low degree test}) If $\tvar_{Q,w} = \tvar_{Q,\ol{w}} =
      v \in \{\alice,\bob\}$, and if $(\tvar_{\Pi,w}, \tvar_{\Pi, \ol{w}}) =
      (\Point_{v}, \ALine_{v})$ (resp.
      $\DLine_v$ instead of $\ALine_v$), execute $\decider^\ld_\ldparams$ on
      input $(\Point, x_{\Pi,w} ,\ALine, x_{\Pi,\overline{w}}, a_w, a_{\ol{w}})$
      (resp.\ $\DLine$ instead of $\ALine$), where $\ldparams = (q,m,d,1)$.
      Reject if $\decider^\ld_{\ldparams}$ rejects.
      \label{enu:ar-input-ld}

    \item (\textbf{Proof encoding checks}): If $\tvar_{Q,w} = \tvar_{Q,\ol{w}} =
      \oracle$,
      \begin{enumerate}
			\item (Consistency test) If $(\tvar_{\Pi,w}, \tvar_{\Pi, \ol{w}}) =
        (\Point_i, \Point_6)$ for some $i \in \{3, \dots, 5\}$, reject if
        $\alpha_i \neq \alpha'_i$.
        \label{enu:ar-proof-cons-ld}

      \item (Individual low degree test) If $(\tvar_{\Pi,w}, \tvar_{\Pi,
          \ol{w}}) = (\Point_i, \ALine_i)$ (resp.\ $\DLine_i$) for some $i \in
        \{3, \dots, 5\}$, execute $\decider^\ld_{\ldparams}$ on input $(\Point,
        x_{\Pi,w} ,\ALine,x_{\Pi,\overline{w}}, a_w,a_{\ol{w}})$ (resp.\ $\DLine$).
        Reject if $\decider^\ld_{\ldparams}$ rejects.
        \label{enu:ar-proof-id-ld}

      \item (Simultaneous low degree test) If $(\tvar_{\Pi, w}, \tvar_{\Pi,
          \ol{w}}) = (\Point_6, \ALine_6)$ (resp.\ $\DLine_6$), execute
        $\decider^\ld_{\ldparams'}$ on input $(\Point, x_{\Pi,w}, \ALine,
        x_{\Pi,\overline{w}}, a_w, a_{\ol{w}})$ (resp.\ $\DLine$), where
        $\ldparams'=(q, m', d, m'+6)$.
        Reject if $\decider^\ld_{\ldparams'}$ rejects.
        \label{enu:ar-proof-sim-ld}
      \end{enumerate}
      \label{enu:ar-proof-encoding}

    \item (\textbf{Game check}): If $\tvar_{Q,w} = \oracle$, then for $v \in
      \{\alice,\bob\}$, compute $x_{w,v} = L^{v}(x_{Q,w})$.
      If $\tvar_{\Pi,w} = \Point_6$, reject if $\pcpverifier((\decider, n , T,
      \qlen, \gamma, x_{w,\alice}, x_{w,\bob}), (z, a_{w} ))$ rejects.
      Otherwise, accept. 
      \label{enu:ar-game}
    \end{enumerate}
  }
  \end{gamespec}
  
  \caption{The decision procedure $\tdans$.
    Parameters $T,Q,q,m,m',d,\gamma$ are defined in
    Section~\ref{sec:ar-params}.}
  \label{fig:decider-pcp}
  
\end{figure}

\subsubsection{Main theorem for answer reduction}

\begin{theorem}\label{thm:ar}
  There exists a polynomial time Turing machine $\ComputeAnsVerifier$ that, on
  input $(\verifier, \lambda, \mu, \gamma)$ where $\verifier$ is an $\ell$-level
  normal form verifier and $\lambda, \mu, \gamma \in \N$, outputs the
  description of the answer reduced verifier $\vans = (\sans, \dans)$.
  Let $T(n)$ and $Q(n)$ be the functions specified in \cref{eq:ar-params-1}, and
  suppose that $\verifier$ satisfies the complexity conditions specified in
  \cref{eq:ar-time-assumption}.
  Then $\vans$ is $\max\{\ell+2, 5\}$-level, has time complexity bounds
  \begin{align*}
    \TIME_{\sans}(n)
    & = O \bigl( \TIME_{\tsora}(n) + \poly(\log(T(n)),|\decider|,\gamma) \bigr) 
     \,=\, \poly\bigl( (\lambda n)^\mu, |\decider|,\gamma \bigr)\;,\\
    \TIME_{\dans}(n)
    & = \poly \bigl( \log(T(n)),Q(n),|\decider|,\gamma \bigr) \,=\,
      \poly\bigl( (\lambda n)^\mu, |\decider|,\gamma \bigr)\;,
  \end{align*}
  and furthermore the sampler $\sampler^\ar$ only depends on
  $\sampler$, the parameter tuple $(\lambda, \mu, \gamma)$, and the description
  length $|\decider|$ (but nothing else about $\decider$).
	Moreover, there exists a function
  \begin{equation*}
    \delta(\eps, n) = a \gamma^a \bigl( (\lambda \cdot |\decider| \cdot n)^{a \mu}
    \cdot \eps^b + (\lambda \cdot |\decider| \cdot n)^{-\mu b\gamma} \bigr)
  \end{equation*}
  for universal constants $a>0$, $0<b<1$, such that the following hold for all $n \geq 2$:
  \begin{enumerate}
  \item (\textbf{Completeness}) If $\verifier_n$ has a projective, consistent,
    and commuting (PCC) strategy of value $1$, then $\vans_n$
    has an SPCC strategy with value $1$.
    \label{enu:ar-completeness}
	\item (\textbf{Soundness}) If $\val^*(\vans_n) > 1 - \eps$ for some $\eps > 0$
    then $\val^*(\verifier_n) \geq 1 - \delta(\eps, n)$.\label{enu:ar-soundness}
  \item (\textbf{Entanglement}) Let $\Ent(\cdot)$ be as defined in \cref{def:ent}.
    Then for all $\eps\geq 0$,
    \begin{equation*}
      \Ent(\vans_n, 1 - \eps) \geq \Ent(\verifier_{n}, 1 - \delta(\eps, n))\;.
    \end{equation*}
  \end{enumerate}
\end{theorem}
Observe that $\TIME_{\dans}(n)$ is polynomial in the \emph{logarithm}
of $T(n)$, the runtime of the decider of $\verifier$, achieving the
desired reduction in runtime and answer complexity.

\begin{proof}
The Turing machine $\ComputeAnsVerifier$ can be described as follows: 
\begin{enumerate}
\item From the description of $\verifier$, compute the description of the
  (typed) oracularized verifier $\tvora = (\tsora,\tdora)$ using the Turing
  machine $\ComputeOracleVerifier$ from \Cref{thm:oracle-completeness}.
\item From the description of the typed sampler $\tsora$ and the parameter tuple
  $(\lambda,\mu,\gamma)$, compute the descriptions of the typed sampler $\tsans$
  as described in \Cref{sec:ar-verifier}.
\item From the descriptions of the decider $\decider$ and the parameter tuple
  $(\lambda,\mu,\gamma)$, compute the descriptions of the typed decider $\tdans$
  as described in \Cref{fig:decider-pcp}.
\item Compute the detyped verifier $\verifier^\ar = (\sampler^\ar,\decider^\ar)
  = \detype(\tvans)$.
\end{enumerate}
The Turing machine $\ComputeAnsVerifier$ takes time that is
$\poly(|\verifier|,\log \lambda,\log \mu,\log \gamma)$, because each step,
including the detyping procedure, runs in time that is polynomial in the length
of the input $(\verifier,\lambda,\mu,\gamma)$.

\paragraph{Properties of the sampler $\sans$.}
It can be verified via inspection that the typed sampler $\tsans$ depends on
$\tsora$ (which itself only depends on $\sampler$; see
Theorem~\ref{thm:oracle-completeness}), and $\tspcp$ (which only depends on the
parameter tuple $(\lambda,\mu,\gamma)$ and the description length $|\decider|$).
Using Theorem~\ref{thm:oracle-completeness}, $\tsora$ is a $\ell$-level typed
sampler and, therefore, $\tsans$ is a $\max\{\ell, 3\}$-level sampler as
$\tspcp$ is a typed $3$-level sampler.
Using Lemma~\ref{lem:detyping-verifiers} for the detyping it follows that
$\sans$ is a $\max\{\ell+2, 5\}$-level typed sampler.
  
\paragraph{Complexity.}
In addition to the running time of $\tsora$, the time complexity of $\tsans$
also includes the running time of $\tspcp$, which is dominated by the complexity
of computing the CL functions $L_{\Point_i}$, $L_{\ALine_i}$, and $L_{\DLine_i}$.
From \Cref{lem:ld-complexity}, this takes time 
\[ \poly(m', \log q) =
\poly(\gamma,s) = \poly(\log T, \log n, Q, |\decider| ,\gamma) = \poly(\log
T,|\decider|,\gamma)\;,\]
where the last equality uses the formulas~\eqref{eq:ar-params-1}. Therefore the time complexity of $\tsans$ is
\[
	O \left ( \TIME_{\tsora}(n) + \poly(\log T(n),|\decider|,\gamma) \right)\;.
\] 
The complexity of $\sans = \detype(\tsans)$ then follows by
Lemma~\ref{lem:detyping-verifiers} and \Cref{eq:ar-time-assumption}.



The decider $\tdans$ executes subroutines $\decider^\ld$ and $\pcpverifier$.
The runtime of $\decider^\ld$ is $\poly(m, d, m', \log q)$.
The runtime of $\pcpverifier$ is given in Theorem~\ref{thm:pcp-decider}; for $T$
and $Q$ as in~\eqref{eq:ar-params-1} it is $\poly(\log T,\log n,\log q, Q,
|\decider|)$.
  
In addition to theese subroutines, $\tdans$ computes the CL functions $L^\alice$
and $L^\bob$ of the oracularized sampler $\tsora$ in Step~\ref{enu:ar-game} of
\Cref{fig:decider-pcp}.
By \Cref{thm:oracle-completeness}, the complexity of $\tsora$ is polynomial in
the time complexity of $\sampler$, which by \Cref{eq:ar-time-assumption} is at
most $Q(n)$.
  
Thus the overall time complexity of $\tdans$, using both the PCP parameter
settings of \Cref{def:pcpparams} and the assumptions of
\Cref{eq:ar-time-assumption}, is $\poly( (\lambda n)^\mu, |\decider|, \gamma)$.
The complexity of $\detype(\tdans)$ follows by
Lemma~\ref{lem:detyping-verifiers}.







\paragraph{Completeness, Soundness, and Entanglement.}
The Completeness property is established in \cref{sec:ar-completeness} and the
Soundness and Entanglement properties are proven in \cref{sec:ar-soundness}.
\end{proof}

\subsection{Completeness of the answer-reduced verifier}
\label{sec:ar-completeness}

We establish the Completeness property of the answer reduced verifier as stated
in \cref{thm:ar}.



\begin{proof}[Proof of the Completeness part of \cref{thm:ar}]
We establish the Completeness property of $\tvans$, which by the Detyping
  Lemma (\Cref{lem:detyping-verifiers}) establishes the Completeness property of
  $\vans$.

  Let $n\geq 1$ be an index for $\tvans$.
  Let $\strategy$ be a PCC strategy for $\verifier_n$ with value $1$.
  By Theorem~\ref{thm:oracle-completeness} it follows that there exists a
  symmetric PCC strategy $\strategy^\ora$ using a state $\ket{\psi}$ with value
  $1$ for $\tvora_n$.
  We define a strategy $\strategy^\ar$ for the typed verifier $\tvans_n$ as
  follows.
  The shared state is $\ket{\psi}$.
  Given the index $n$ and $(\lambda,\mu,\gamma)$, each player can compute
  $(q,m,d,m',s) = \pcpparams(n, T, \qlen,\sigma,\gamma)$
  (see~\Cref{def:pcpparams}), where $T=T(n)$ and $Q=Q(n)$ are as
  in~\eqref{eq:ar-params-1}, and $\sigma = \abs{\decider}$.
  Let $M = 2^m$.

  \begin{enumerate}
  \item On receipt of a question $((\tvar_Q,\tvar_\Pi),(x_Q, x_\Pi))$ a player
    first measures their share of $\ket{\psi}$ using the projective measurement
    for $\strategy^\ora$ for the typed question $(\tvar_Q,x_Q)$ to obtain an
    outcome $a_Q$.
    The player then computes an answer, depending on $\tvar_Q, a_Q$ and
    $\tvar_\Pi,x_\Pi$, as follows:
    \begin{enumerate}
    \item Suppose $\tvar_Q = v \in \{\alice,\bob\}$ and $\tvar_\Pi \in
      \{\Point_v, \ALine_v \}$ (resp. $\DLine_v$ instead of $\ALine_v$).
      Let $a_Q' = a_Q$ if $a_Q$ has length at most~$T$, and let $a_Q'$ be the
      truncation of~$a_Q$ to its first~$T$ symbols otherwise.
      Let $\ell_a \leq T$ be the length of~$a_Q'$, and set 
      \begin{equation*}
      a_Q'' = \mathrm{enc}_\Gamma(a_Q', \sqcup^{M/2 - \ell_a}).
      \end{equation*}
      Next, the player computes the low-degree encoding $g_{a_Q''}$ of $a_Q''$
      using the low-degree encoding described in \Cref{sec:ld-encoding}, in the
      same way as described in Section~\ref{sec:ar-pcp}.
      The player then returns the restriction of $g_{a_Q''}$ to the
      axis-parallel line (resp.
      the diagonal line) specified by $x_\Pi$.
\item
If $\tvar_Q=\oracle$, for $v \in \{\alice, \bob\}$ the player computes
      questions $x_v = L^{v}(x_Q)$, as in Step \ref{enu:oracle-game} of
      $\tdora$.
      The player parses $a_Q$ as a pair $(a_\alice, a_\bob)$.
      Let $a_\alice' = a_\alice$ if $a_\alice$ has length at most~$T$, and let
      $a_\alice'$ be the truncation of~$a_\alice$ to its first~$T$ symbols
      otherwise.
      Let $\ell_{\alice} \leq T$ be the length of~$a_\alice'$, and set
      \begin{equation*}
        a_\alice'' = \mathrm{enc}_\Gamma(a_\alice', \sqcup^{M/2 - \ell_{\alice}}).
      \end{equation*}
      Define $a_\bob''$ similarly.
      The player computes a PCP proof $\Pi=(g_1, \dots, g_5, c_0, \dots, c_{m'})$ as
      described in the completeness case of \Cref{thm:pcp-decider} for the tuple
      $(\decider, n, T, x_\alice, x_\bob)$, where the polynomials $g_1, g_2$ are
      low-degree encodings of $a_\alice''$ and $a_\bob''$, respectively.
      \begin{enumerate}
      \item If $\tvar_\Pi \in \{\Point_i, \ALine_i\}$ (resp.
        $\DLine_i$ instead of $\ALine_i$) for $i \in \{1, \dots, 5\}$, the
        player returns the restriction of $g_i$ to the axis-parallel line (resp.
        diagonal line) of $\F_q^{m}$ specified by $x_\Pi$.
      \item If $\tvar_\Pi \in \{\Point_6, \ALine_6\}$ (resp.
        $\DLine_6$ instead of $\ALine_6$), the player returns the restriction of
        all the polynomials $g_1, \dots, g_5,$ $c_0, \dots, c_{m'}$ to the
        axis-parallel line of $\F_q^{m'}$ (resp.
        the diagonal line) specified by $x_\Pi$.
      \end{enumerate}
    \item In all other cases the player returns $0$.
    \end{enumerate}
  \end{enumerate}
  The strategy $\strategy^\ar$ is projective and consistent because $\strategy^\ora$
  is.
  To show that it has value $1$, we first observe that by definition it
  satisfies all consistency checks.
  Moreover, the strategy passes all low-degree tests with certainty because it
  always returns restrictions of consistent polynomials.
  Finally, it also passes the game check with probability $1$.
  This follows from the completeness statement of the PCP made in
  Theorem~\ref{thm:pcp-decider} and the fact that, if $\decider$ accepts the
  input $(n, x_\alice,x_\bob, a_\alice,a_\bob)$ in time at most $T$ then it also
  accepts $(n, x_\alice,x_\bob, a'_\alice,a'_\bob)$ in time at most $T$, where
  $a'_\alice$ and $a'_\bob$ are obtained from $a_\alice$ and $a_\bob$ by
  truncating them to strings of length~$T$ if their lengths exceed~$T$.

	To show that $\strategy^\ar$ is commuting, note that using the product structure
  of $\tsans$ every typed question pair with positive probability consists
  of a pair of questions $((\tvar_{Q,\alice},x_{Q, \alice}),(\tvar_{Q,\bob},
  x_{Q,\bob})) $ with positive probability for $\tsora$, together with an
  arbitrary pair $((\tvar_{\Pi,\alice},x_{\Pi, \alice}),
  (\tvar_{\Pi,\bob}, x_{\Pi,\bob}))$.
  Using that $\strategy^\ora$ is commuting and that the additional operations
  associated with $((\tvar_{\Pi,\alice},x_{\Pi, \alice}),
  (\tvar_{\Pi,\bob}, x_{\Pi,\bob}))$ amount to classical post-processing it
  follows that $\strategy^\ar$ is commuting.

  This establishes the existence of a symmetric PCC strategy for
  $\tvans_n$ with value $1$.
	By Lemma~\ref{lem:detyping-verifiers} it follows that there exists a symmetric
  PCC strategy for $\vans_n = \detype(\tvans)_n$ with value $1$.

\qedhere

\end{proof}

\subsection{Soundness of the answer-reduced verifier}
\label{sec:ar-soundness}



\begin{proof}[Proof of the soundness part of \cref{thm:ar}]
  We first show the soundness for the typed verifier $\tvans$.
  Soundness for the detyped verifier $\vans$ follows from
  Lemma~\ref{lem:detyping-verifiers}, with a constant-factor loss using that the
  type set $\type^\ar$ for $\tvans$ has constant size.
	
  We proceed in two steps.
  Fix an index $n\geq 1$ and suppose that $\val^*(\tvans_n) > 1-\eps$ for
  some $\eps > 0$.
  Observe that $\tsora$ and $\tspcp$ both sample distributions that
  are invariant under permutation of the two players; therefore, the same holds
  for $\tsans$.
  Moreover, the decider $\tdans$ treats both players symmetrically.
  Therefore, the game played by $\tvans_n$ is a symmetric game.
  Applying Lemma~\ref{lem:symmetric-strat} it follows that there exists a
  symmetric projective strategy $\strategy = (\ket{\psi},M)$ for
  $\tvans_n$ with value greater than $1 - \eps$.
	
  We use the following shorthand notation.
  A pair of questions to the players is
  $((\tvar_\alice,x'_\alice),(\tvar_\bob,x'_\bob))$ where for
  $w\in\{\alice,\bob\}$, $\tvar_w=(\tvar_{Q,w},\tvar_{\Pi,w})$ and
  $x'_w=(x'_{Q,w}, x'_{\Pi,w})$.
  When $w$ is clear from context we omit it from the subscript.
  Fixing a $w$, whenever $\tvar_Q = \oracle$ we introduce $x_\alice =
  L^\alice(x_Q)$ and $x_\bob =L^\bob(x'_Q)$ and often write directly the player's
  question as $x_Q = (x_\alice,x_\bob)$.
  Whenever $\tvar_Q=v \in \{\alice,\bob\}$ we slightly abuse notation and write
  the question as $x_Q = (x_{v},v)$, explicitly including the type to clarify
  which player it points to.
	
  We denote the measurements used by both players in strategy $\strategy$ by
  $\{M({x}_Q)^{ {x}_{\Pi}}_{a}\}$, where for the sake of clarity we have
  notationally separated the two parts ${x}_Q$ and $x_\Pi$ of the question and
  omitted explicit mention of the associated types $\tvar_Q$ and $\tvar_\Pi$ (we
  include the type and write $M({x}_Q)^{\tvar_\Pi, {x}_{\Pi}}_{a}$ when it is
  needed for clarity).
  First we show that the strategy $\strategy$ is close to a strategy
  $\strategy'$ that performs ``low-degree'' measurements: upon receipt of a
  typed question $(\tvar,x)=((\tvar_Q,\tvar_\Pi),({x}_Q,x_\Pi))$ a player first
  performs a measurement depending on ${x}_Q$ to obtain a tuple of low-degree
  polynomials, and then returns evaluations of those polynomials on the
  subspaces (either axis-parallel or diagonal lines) specified by $x_\Pi$.
  This step of the argument uses the quantum soundness of the low-degree test
  performed in Steps~\ref{enu:ar-input-ld}
  and~\ref{enu:ar-proof-encoding} of Figure~\ref{fig:decider-pcp}.
  Next, we ``decode'' this strategy to produce a strategy $\strategy''$ for
  $\tvora$ with a high value.
  This step makes use of the classical soundness of the underlying PCP shown in
  \Cref{sec:pcp-cktval-new}.
  The conclusion of the theorem then follows from the soundness of
  $\tvora$ (\Cref{thm:oracle-soundness}).
  We proceed with the details.

	We start by showing a sequence of claims that establish approximations implied
  by the assumption that $\strategy^\ar$ succeeds with probability greater than 
  $1-\eps$ in the decision procedure implemented by the decider in
  Figure~\ref{fig:decider-pcp}.


  \begin{claim}[Global consistency check, Step~\ref{enu:ar-global-consistency}]
    \label{claim:ar-1}
    On average over questions
    $(\tvar_\alice,x_\alice)=((\tvar_Q,\tvar_\Pi),({x}_Q, x_{\Pi}))$ sampled
    from the marginal distribution of $ \mu_{\tsans}$ on the first player
    it holds that
    \begin{equation}
      \label{eq:global-consistency}
      M({x}_Q)^{x_{\Pi}}_a \ot I \simeq_{\eps} I \ot M({x}_Q)^{x_{\Pi}}_a\;.
    \end{equation}
  \end{claim}

  \begin{proof}
    First we observe that the condition $\tvar_\alice = \tvar_\bob$ for the
    global consistency check, Step~\ref{enu:ar-global-consistency} in
    Figure~\ref{fig:decider-pcp}, holds with constant probability over the
    choice of a pair of questions $(\tvar_\alice,x_\alice),(\tvar_\bob,x_\bob)$
    sampled according to $\mu_{\tsans}$.
    Thus $\strategy$ must succeed in this test with probability $1-O(\eps)$,
    conditioned on the test being executed: this is because each of
    $\tsora$ and $\tspcp$ have a constant probability of returning
    a pair of questions of the same type.

    Moreover, observe that conditioned on $\tvar_\alice = \tvar_\bob$ a pair of
    questions $((\tvar_\alice,x_\alice),(\tvar_\alice,x_\bob))\sim
    \mu_{\tsans}$ is such that $x_\alice = x_\bob = L_{\tvar_\alice}(z)$,
    where $z$ is the sampler seed and $L_{\tvar_\alice}$ the CL function of type
    $\tvar_\alice$ associated with $\tsans$.
    The claim then follows directly from the test and the definition of
    approximate consistency (Definition~\ref{def:consistency}).
  \end{proof}

  \begin{claim}[Input consistency check, Step~\ref{enu:ar-input-consistency}]
    \label{claim:ar-2}
    For all $v \in \{\alice,\bob\}$, on average over question pairs
    $({x}_\alice, {x}_\bob) \sim \mu_{\sampler}$ and $z=(y_1,\ldots,y_5,o,w)\in
    \F_q^{m'}$ sampled uniformly at random,
    \begin{equation}
      M({x}_\alice, {x}_\bob)^{\Point_6, z}_{\alpha_v} \ot I \simeq_{\eps} I \ot
      M(x_v,v)^{ \Point_v, y_v}_{\alpha_v}\;, \label{eq:input-consistency}
    \end{equation}
    where as in Remark~\ref{rk:ab-01} we made the identification
    $1\leftrightarrow \alice$ and $2\leftrightarrow \bob$.
    Moreover, an analogous relation holds for operators acting on opposite sides
    of the tensor product.
  \end{claim}

	\begin{proof}
    For $w=\alice$ and fixed $v\in\{\alice,\bob\}$ there is a constant
    probability that $\tvar_{Q,w} = \oracle$, $\tvar_{Q,\ol{w}}=v$, and
    $(\tvar_{\Pi,w},\tvar_{\Pi, \ol{w}}) = (\Point_6, \Point_{v})$.
    Therefore, the input consistency check in Step~\ref{enu:ar-input-consistency} is
    executed with constant probability, and $\strategy$ must pass it with
    probability $1-O(\eps)$, conditioned on the test being executed.

    Moreover, conditioned on $\tvar_{Q,w} = \oracle$, $\tvar_{Q,\ol{w}}=v$, and
    $(\tvar_{\Pi,w},\tvar_{\Pi, \ol{w}}) = (\Point_6, \Point_{v})$, the
    distribution of $(x_{Q,w},x_{Q,\overline{w}})$ is $((x_\alice,x_\bob),x_v)$
    for $(x_\alice,x_\bob)\sim \mu_\sampler$ and the distribution of
    $(x_{\Pi,w},x_{\Pi,\overline{w}})$ is $(z,y_v)$ for a uniformly random $z
    \in \F_q^{m'}$.
    Eq.~\eqref{eq:input-consistency} then follows directly from the
    specification of the test and the definition of approximate consistency.
    The ``moreover'' part follows from the case $w=\bob$.
  \end{proof}
	
	\begin{claim}[Input low degree test, Step~\ref{enu:ar-input-ld}]\label{claim:ar-3}
    For each $v \in \{\alice,\bob\}$ and for each $x$ in the support of the
    marginal of $\mu_S$ on player $v$ there exists a measurement
    $\{G^{{x},v}_{g}\} \in \polymeas{m}{d}{q}$ such that the following hold for
    some $\delta_1 = O(\delta_{\ld}(O({\eps}), q, m, d, 1))$, where
    $\delta_{\ld}$ is defined in Theorem~\ref{lem:ld-soundness}.
    For all $v\in\{\alice,\bob\}$, on average over ${x}$ chosen from the
    marginal of $\mu_{\sampler}$ on player $v$ and $y_v\in \F_q^m$ sampled
    uniformly at random,
    \begin{align}
      M({x}, v)^{\Point_v,y_v}_{\alpha}\ot I
      & \simeq_{\delta_1} I \ot G^{{x}, v}_{[\eval_{y_v}(\cdot) =
        \alpha]}\;, \label{eq:ld-single-input} \\
      I\ot  M({x}, v)^{\Point_v,y_v}_{\alpha}
      & \simeq_{\delta_1}  G^{{x}, v}_{[\eval_{y_v}(\cdot) =
        \alpha]} \ot I\;, \label{eq:ld-single-input-b} \\
      G^{{x}, v}_{g} \ot I
      & \simeq_{\delta_1} I \ot G^{{x},
        v}_{g}\;,\label{eq:ld-self-input}
    \end{align}
    where we used the notation $\eval_{y_v}(g)=g(y_v)$ for the evaluation map.
	\end{claim}
	
	\begin{proof}
    Fix $v\in\{\alice,\bob\}$.
    For any $w\in\{\alice,\bob\}$ there is a constant probability that
    $\tvar_{Q,w} = \tvar_{Q,\ol{w}} = v$ and $(\tvar_{\Pi,w}, \tvar_{\Pi,
      \ol{w}}) = (\Point_{v}, \ALine_{v})$ (resp.\ $\DLine_v$).
    Therefore, the input low degree test in Step~\ref{enu:ar-input-ld} of
    \Cref{fig:decider-pcp} is executed with constant probability, and
    $\strategy$ must pass it with probability $1-O(\eps)$, conditioned on the
    test being executed.
			
    Observe that by definition the distribution of $(x_{\Pi,\alice},
    x_{\Pi,\bob})$ conditioned on $\tvar_{Q,w} = \tvar_{Q,\ol{w}} = v$,
    uniformly random $x_Q = (x_v,v)$, and $(\tvar_{\Pi,w}, \tvar_{\Pi, \ol{w}})
    = (\Point_{v}, \ALine_{v})$ (resp.
    $\DLine_v$), where $w\in\{\alice,\bob\}$ is uniformly random, is exactly the
    distribution of questions in the game $\game^\ld$ described in
    Section~\ref{sec:ld-game}, parametrized by $\ldparams = (q, m, d, 1)$.
			
    For every $v\in\{\alice,\bob\}$ and question $x=L^v(z)$ in the support of
    the marginal distribution of $\mu_\sampler$ on player $v$ let $\eps_{x,v}$
    be the probability that $\strategy$ is rejected in Step~\ref{enu:ar-input-ld},
    conditioned on the test being executed and on average over
    $w\in\{\alice,\bob\}$.
    Then $\E[\eps_{x,v}] = O(\eps)$, where the expectation is taken over a
    uniformly random $v\in\{\alice,\bob\}$ and $x=L^v(z)$ for uniformly random
    $z$.

    By definition it follows that the strategy $\tsans$ conditioned on the
    first part of the players' questions being $\tvar_{Q,w} = \tvar_{Q,\ol{w}} =
    v$ and $x_{Q,\alice}=x_{Q,\bob}=x$ is a projective strategy that succeeds
    with probability $1-\eps_{x,v}$ in the low-degree test
    $\decider^\ld_\ldparams$ executed in Item~\ref{enu:ar-input-ld}.
	
    We may thus apply Theorem~\ref{lem:ld-soundness} to obtain $\{G^{{x},v}_{g}\}
    \in \polymeas{m}{d}{q}$ such
    that~\eqref{eq:ld-single-input},~\eqref{eq:ld-single-input-b}
    and~\eqref{eq:ld-self-input} each hold with approximation error
    $O(\delta_{\ld}(\eps_{x,v},q,m,d,1))$.
    Using that for fixed $q,m,d$ the function
    $\eps\mapsto\delta_{\ld}(\eps,q,m,d,1)$ is concave, the claim follows from
    Jensen's inequality.
  \end{proof}

  \begin{claim}[Proof encoding checks, Step~\ref{enu:ar-proof-encoding}]
    \label{claim:ar-4}
    For each $x_Q=(x_\alice,x_\bob)$ in the support of $\mu_\sampler$ there
    exist measurements $\{G^{(x_\alice,x_\bob),i}_{g}\} \in \polymeas{m}{d}{q}$
    for each $i\in\{3,4,5\}$ and
    \begin{equation*}\{J^{({x}_\alice, {x}_\bob)}_{f_1, \dots, f_5, c_0,
      \dots, c_{m'}}\} \in \simulpolymeas{m'}{d}{q}{m'+6}
      \end{equation*}
      such that
    the following hold for some
    \begin{equation*}
      \delta_2 = O\big(\delta_{\ld}(O({\eps}), q, m, d, 1)+
      \delta_{\ld}(O({\eps}), q, m', d, m'+6)\big)\;.
    \end{equation*}
    First, for all $i\in \{3,4,5\}$, on average over $({x}_\alice, {x}_\bob)
    \sim \mu_{\sampler}$ and $z=(y_1,\ldots,y_5,o,w)$ of type $\Point_6$ sampled
    uniformly at random,
    \begin{equation}
      I \ot M({x}_\alice, {x}_\bob)^{ \Point_i, y_i}_{\alpha_i} \simeq_{\eps}
      M({x}_\alice, {x}_\bob)^{\Point_6, z}_{\alpha_i} \ot I\;.
      \label{eq:ld-cons-pf}
    \end{equation}
    Second, for all $i\in \{3,4,5\}$ and on average over $({x}_\alice, {x}_\bob)
    \sim \mu_{\sampler}$ and $y_i\in\F_q^m$ sampled uniformly at random,
    \begin{align}
      M({x}_\alice, {x}_\bob)^{\Point_i,  y_i}_{\alpha}\ot I
      &\simeq_{\delta_2} I \ot G^{({x}_\alice, {x}_\bob), i}_{[\eval_{y_i}(\cdot) =
        \alpha]}\;, \label{eq:ld-single-pf} \\
      G^{({x}_\alice, {x}_\bob), i}_{g} \ot I
      &\simeq_{\delta_2} I \ot G^{({x}_\alice, {x}_\bob),
        i}_{g}\;. \label{eq:ld-self-pf}
    \end{align}
    Third, for all $i\in\{1,\ldots,5\}$ and $j\in\{0,\ldots,m'\}$, on average
    over $({x}_\alice, {x}_\bob) \sim \mu_{\sampler}$ and $z\in\F_q^{m'}$
    sampled uniformly at random,
    \begin{align}
      M({x}_\alice, {x}_\bob)^{\Point_6,  z}_{\alpha_1, \ldots, \alpha_5,
      \beta_0, \ldots, \beta_{m'}} \ot I
      & \simeq_{\delta_2} I \ot J^{({x}_\alice, {x}_\bob) }_{[\eval_{z}(\cdot) =
        (\alpha_1,\ldots,\alpha_5,\beta_0,\ldots,\beta_{m'})]}\;,
        \label{eq:ld-simul-a}\\
      J^{({x}_\alice,{x}_\bob)}_{f_1, \dots, f_5, c_0, \dots, c_{m'}} \ot I
      &\simeq_{\delta_2} I \ot J^{({x}_\alice, {x}_\bob)}_{f_1, \dots, f_5,
        c_0, \dots, c_{m'}}\;.\label{eq:ld-self-simul}
    \end{align}
    Moreover, analogous equations
    to~\eqref{eq:ld-cons-pf},~\eqref{eq:ld-single-pf} and~\eqref{eq:ld-simul-a}
    hold with operators acting on opposite sides of the tensor product.
    Here, recall the definition of $\polymeas{\cdot}{\cdot}{\cdot}$ from
    \Cref{def:ld-meas}.

\end{claim}
	
	\begin{proof} 
    The proof of the first item is similar to the proof of
    Claim~\ref{claim:ar-2}, and we omit it.

    The proof of the second and third items is similar to the proof of
    Claim~\ref{claim:ar-3}, and we include more details.
    Fix an $i\in\{3,4,5\}$.
    For any $w\in\{\alice,\bob\}$ there is a constant probability that
    $\tvar_{Q,w} = \tvar_{Q,\ol{w}} = \oracle$ and $(\tvar_{\Pi,w}, \tvar_{\Pi,
      \ol{w}}) = (\Point_i, \ALine_i)$ (or $\DLine_i$ instead of $\ALine_i$), in
    which case the individual low-degree test in~Step~\ref{enu:ar-proof-id-ld}
    is executed.
    Therefore, $\strategy$ must succeed in that part of the test with
    probability $1-O(\eps)$ conditioned on the test being executed.
					
    Furthermore, for fixed $i\in \{3,4,5\}$ and uniformly random
    $w\in\{\alice,\bob\}$, conditioned on the test being executed for that $i$
    and $w$ the distribution of $(x_{\Pi,\alice},x_{\Pi,\bob})$ is exactly the
    distribution of questions in the game $\game^\ld$ described in
    Section~\ref{sec:ld-game}, parametrized by $\ldparams = (q, m, d, 1)$.
			
    For every $i\in\{3,4,5\}$ and $x=(x_\alice,x_\bob)$ in the support of
    $\mu_\sampler$ let $\eps_{x,i}$ be the probability that $\strategy$ is
    rejected in~Step~\ref{enu:ar-proof-id-ld}, conditioned on the test being
    executed for that $i$ and on average over $w\in\{\alice,\bob\}$.
    Then for each $i$, $\E[\eps_{x,i}] = O(\eps)$, where the expectation is
    taken over a uniformly random $x\sim \mu_\sampler$.

    By definition of the individual low-degree test it follows
    from~\Cref{lem:ld-soundness} that for every $x=(x_\alice,x_\bob)$ in the
    support of $\mu_\sampler$ and $i\in\{3,4,5\}$ there is a measurement
    $\{G^{({x}_\alice, {x}_\bob), i}_{g}\} \in \polymeas{m}{d}{q}$ such that on
    average over $y_i\in \F_q^{m'}$ of type $\Point_i$ sampled uniformly at
    random, Eq.~\eqref{eq:ld-single-pf} and~\eqref{eq:ld-self-pf} both hold with
    approximation error $O(\delta_{\ld}(\eps_{x,i},q,m,d,1))$.
    Eq~\eqref{eq:ld-single-pf} and~\eqref{eq:ld-self-pf} follow using the concavity
    of $\delta_{\ld}$ as a function of $\eps$.

    Finally we consider the simultaneous low-degree
    test,~Step~\ref{enu:ar-proof-sim-ld}.
    Here as well, using that there is a constant probability that $\tvar_{Q,w} =
    \tvar_{Q,\ol{w}} = \oracle$ and $(\tvar_{\Pi,w}, \tvar_{\Pi, \ol{w}}) =
    (\Point_6, \ALine_6)$ (resp.
    $\DLine_6$) it follows that $\strategy^\ar$ must succeed in that part of the
    test with probability $1-O(\eps)$.
    Using a similar argument as before it follows from~\Cref{lem:ld-soundness}
    (this time for parameters $(q,m',d,m'+6)$) that for every
    $(x_\alice,x_\bob)$ there is a family of measurements $\{J^{({x}_\alice,
      {x}_\bob)}_{f_1, \dots, f_5, c_0, \dots, c_{m'}}\} \in
    \simulpolymeas{m'}{d}{q}{m'+6}$ such that on average over $z\in \F_q^{m'}$
    sampled uniformly at random, Eq.~\eqref{eq:ld-simul-a}
    and~\eqref{eq:ld-self-simul} both hold with approximation error
    $O(\delta_{\ld}(O({\eps}), q, m', d, m'+6))$.
  \end{proof}

  The families of measurements $\{G^{x_Q,i}_g\}$ and
  $\{J^{x_Q}_{f_1,\ldots,f_5,c_0,\ldots,c_{m'}}\}$, for $x_Q$ in the support of
  $\mu_\sampler$ and $i\in\{1,\ldots,5\}$, whose existence follows from
  Claim~\ref{claim:ar-3} and Claim~\ref{claim:ar-4} have outcomes that are low
  (individual) degree polynomials: for the first family, individual degree $d$
  polynomials $g:\F_q^m\to \F_q$, and for the second, tuples of individual
  degree $d$ polynomials $f_i,c_j:\F_q^{m'} \to \F_q$.
  Recall that $m' = 5m + 5 +s$ and that an element $z\in \F_q^{m'}$ is written
  as a triple $(y,o,w)$ with $x=(y_1,\ldots,y_5)\in \F_q^{5m}$, $o\in\F_q^5$ and
  $w\in\F_q^s$.
  The following claim, whose proof is based on~\Cref{lem:ld-sandwich}, shows
  that we can reduce to a situation where the polynomials $f_1,\ldots,f_5$
  returned by $J$ are such that for each $i\in \{1,\ldots,5\}$, $f_i$ only
  depends on the $y_i$, and not on the entire variable $z$.
	
	\begin{claim}\label{claim:ar-5}
    For all $(x_\alice,x_\bob)$ in the support of $\mu_S$ and individual degree $d$
    polynomials $g_1,\ldots,g_5:\F_q^m\to \F_q$ and $c_0,\ldots,c_{m'}:\F_q^{m'}
    \to \F_q$ define
    \begin{equation}
      \Lambda^{{x}_\alice, {x}_\bob}_{ g_1, \dots, g_5, c_0, \dots, c_{m'}} =
      G^{{x}_\alice, 1}_{g_1} G^{{x}_\bob, 2}_{g_2} G^{({x}_\alice, {x}_\bob),
        3}_{g_3} G^{({x}_\alice,{x}_\bob),4}_{g_4} G^{({x}_\alice, {x}_\bob),5}_{g_5} J^{({x}_\alice, {x}_\bob)}_{c_0,
        \dots, c_{m'}} G^{({x}_\alice, {x}_\bob), 5}_{g_5}
      G^{({x}_\alice,{x}_\bob),4}_{g_4}
      G^{({x}_\alice, {x}_\bob), 3}_{g_3} G^{x_\bob,2}_{g_2} G^{{x}_\alice,1}_{g_1}\;,
      \label{eq:lambda-giant-sandwich}
    \end{equation}
    where the outcomes $(f_1,\ldots,f_5)$ of the $J$ operator in the middle have
    been marginalized over.
    Then there is a
    \[\delta_3 = O\Big(\Big(\delta_{\ld}({\eps}, q, m', d,
      m'+6)^{1/2}+ \frac{m'd}{q}\Big)^{1/2}\Big)\]
such that
    \[ \delta_3 \,\geq\, \max\big\{ \eps, \,\delta_1,\,\delta_2\big\}\]
    and on average over $({x}_\alice, {x}_\bob) \sim
    \mu_{\sampler}$ and $z\in \F_q^{m'}$ sampled uniformly at random,
    \begin{equation}
      \Lambda^{{x}_\alice, {x}_\bob}_{[\eval_{z}(\cdot) = ( \alpha, \beta)]} \ot I
      \simeq_{\delta_3} I \ot J^{{x}_\alice, {x}_\bob}_{[\eval_{z}(\cdot)
        = ( \alpha, \beta)]}\;.
    \label{eq:lambda-consistent}
  \end{equation}
	Moreover, a similar equation holds with the operators acting on opposite sides
  of the tensor product.

\end{claim}
	\begin{proof}
    We apply~\Cref{lem:ld-sandwich} with the following setting of parameters.
    The number of sets of functions $k$ is set to $6$.
    The question set $\mathcal{X}$ is set to the support of $\mu_\sampler$, and
    the distribution $\mu$ on it is the distribution $\mu_\sampler$.
    The sets $\mathcal{G}_i$ for $i \in \{1, \dots, 5\}$ consist of individual
    degree $d$ polynomials over $\F_q^{m'}$ that depend only on the $i$-th block
    of $m$ variables.
    The set $\mathcal{G}_6$ consists of $(m' + 1)$-tuples of individual degree
    $d$ polynomials over $\F_q^{m'}$.

    We first verify the assumption on the sets of functions.
    Since all polynomials have individual degree at most $d$ (and therefore
    total degree at most $m'd$), by~\Cref{lem:schwartz-zippel} the parameter
    $\eps$ in Lemma~\ref{lem:ld-sandwich} can be set to $m'd/q$.

    The family of measurements $\{A^x_{g_1,\ldots,g_6}\}$ in
    Lemma~\ref{lem:ld-sandwich} is the family of measurements
    $\{J^{(x_\alice,x_\bob)}_{f_1,\ldots,f_5,c_0,\ldots,c_{m'}}\}$ here, where
    we set $g_i = f_i$ for $i\in \{1,\ldots,5\}$ and $g_6 =
    (c_0,\ldots,c_{m'})$.
    The measurements $\{G^{i,x}_g\}$ in Lemma~\ref{lem:ld-sandwich} are
    $\{G^{x_\alice,i}_{g}\}$ for $i\in \{1,2\}$,
    $\{G^{(x_\alice,x_\bob),i}_{g}\}$ for $i\in \{3,4,5\}$, and
    $\{J^{(x_\alice,x_\bob)}_g\}$ for $i=6$.
    To ensure that all polynomials are defined over the same range, we treat
    $g:\F_q^m \to \F_q$ that is an outcome of some $\{G^{i,x}_g\}$ as a
    polynomial $g':\F_q^{m'} \to \F_q$, where the role of the $m$ variables of
    $g$ is taken by the $i$-th block of $m$ variables of $g'$.

    We verify that assumption~\eqref{eq:ld-sandwich-1} in the lemma
    holds with an error $\delta = O((\max\{\eps, \delta_1, \delta_2\})^{1/2})$.
    For $i\in\{1,2\}$ the assumption follows by
    combining~\eqref{eq:input-consistency} and~\eqref{eq:ld-single-input-b}
    with~\eqref{eq:ld-simul-a} and Fact~\ref{fact:data-processing}.
    For $i\in \{3,4,5\}$ we use~\eqref{eq:ld-cons-pf} instead
    of~\eqref{eq:input-consistency} and~\eqref{eq:ld-single-pf} instead of
    ~\eqref{eq:ld-single-input-b}.
    Finally, for $i=6$ we use~\eqref{eq:ld-self-simul} and
    Fact~\ref{fact:data-processing}.
    In these derivations, we use Fact~\ref{fact:triangle-for-simeq}, the
    triangle inequality for ``$\simeq$''.

    The conclusion follows from~\Cref{lem:ld-sandwich}, with an error
    $\delta_3 = O((\delta + m'd/q)^{1/2})$. Observing that
    $\eps,\delta_1,\delta_2 = O(\delta_{\ld}(\eps,q,m',d,m'+6))$, as can be
    verified from the definition of $\delta_{\ld}$ given in
    Theorem~\ref{lem:ld-soundness}, and therefore $\delta =
    O(\delta_{\ld}(\eps, q,m',d,m'+6)^{1/2})$, the claimed bound
    \[ \delta_3 = O\Big(\Big( \delta_{\ld}(\eps, q, m', d, m'+6)^{1/2} +
      \frac{m'd}{q} \Big)^{1/2} \Big) \]
    follows.

\end{proof}

At this point, we have constructed measurements $G$ and $\Lambda$ that return
low-degree polynomials in a similar way as is expected from the honest strategy
in $\tvans_n$, as described in the proof of \cref{thm:ar}.
These measurements can be used to specify a new strategy $\strategy'$ for the
game $\tvans_n$ as follows.
The shared state remains the state $\ket{\psi}$ used in $\strategy$.
For $w\in \{\alice,\bob\}$, upon reception of a question $(\tvar_w,x_w)$ player
$w$ performs the following.
If $\tvar_w = (\tvar_{Q,w},\tvar_{\Pi,w})$ is such that $\tvar_{Q,w} = \oracle$,
the player measures their share of $\ket{\psi}$ using the measurement
$\Lambda^{x_{Q,w}}$ defined in Claim~\ref{claim:ar-5} to obtain a tuple
$(g_1,\ldots,g_5,c_0,\ldots,c_{m'})$ of polynomials.
The player then answers exactly as in the strategy described in the
``completeness'' part of the proof of \cref{thm:ar}.
Similarly, if $\tvar_{Q,w} = v \in \{\alice,\bob\}$ the player first measures
their share of $\ket{\psi}$ using the measurement $G^{x_{Q,w},v}$ from
Lemma~\ref{claim:ar-3} to obtain a polynomial $g$ as outcome; the player then
answers according to the same honest strategy.
	
\begin{lemma}\label{lem:ar-ar}
	The strategy $\strategy'$ succeeds with probability $1-O(\delta_3^{1/2})$ in the
  game $\tvans_n$.
\end{lemma}

\begin{proof}
  First we establish useful consistency relations.
  By combining \Cref{eq:lambda-consistent} and \Cref{eq:ld-simul-a} and applying
  Fact~\ref{fact:data-processing} we obtain that for all $i\in \{1,\ldots,5\}$,
  on average over $({x}_\alice, {x}_\bob) \sim \mu_{\sampler}$ and $z\in
  \F_q^{m'}$ sampled uniformly at random,
  \begin{equation}
    \label{eq:m-lambda-con}
    M({x}_\alice, {x}_\bob)^{\Point_6, z}_{\alpha_i} \ot I
    \simeq_{\delta_3} I \ot  \Lambda^{({x}_\alice, {x}_\bob)}_{[\eval_z(\cdot)_i =
      \alpha_i]}\;,
  \end{equation}
	and a similar equation holds with the operators acting on opposite sides of
  the tensor product.
  Next, combining~\eqref{eq:m-lambda-con} together with
  \Cref{eq:input-consistency} and \Cref{eq:ld-single-input} it follows that for
  each $v \in \{\alice,\bob\}$, on average over $({x}_\alice, {x}_\bob) \sim
  \mu_{\sampler}$ and $z=(y_1,\ldots,y_5,o,w) \in \F_q^{m'}$ sampled uniformly
  at random,
  \begin{equation}\label{eq:m-lambda-con-b}
    G^{{x}_v, v}_{[\eval_{y_v}(\cdot) = \alpha_v]} \ot I \simeq_{\delta_3} I \ot
    \Lambda^{({x}_\alice, {x}_\bob)}_{[\eval_z(\cdot)_v = \alpha_v]}\;.
  \end{equation}
  From the Schwartz-Zippel lemma (\Cref{lem:schwartz-zippel}) it follows that
  the probability that any two distinct individual degree $d$ polynomials $g_v$
  (an outcome of $G^{{x}_v, v}$) and $g'_v$ (an outcome of
  $\Lambda^{({x}_\alice, {x}_\bob)}$) agree at a uniformly random point $y_v\in
  \F_q^m$ is at most $m'd/q$.
  It thus follows from~\eqref{eq:m-lambda-con-b} that for all
  $v\in\{\alice,\bob\}$, on average over $({x}_\alice, {x}_\bob) \sim
  \mu_{\sampler}$ and $y_v \in\F_q^M$ sampled uniformly at random,
  \begin{equation}
    \label{eq:ld-strat-con-def}
    G^{{x}_v, v}_{g_v} \ot I \simeq_{\delta_3+m'd/q} I \ot \Lambda^{({x}_\alice,
      x_\bob)}_{g_v}\;.
  \end{equation}
  We now show that $\strategy'$ is accepted by $\tvans_n$ with high
  probability.
  We bound the probability of succeeding in each subtest.

  First note that the strategy is accepted in Step 1 of \Cref{fig:decider-pcp}.
  For the $G$ measurements, consistency follows from~\eqref{eq:ld-self-input}.
  For the $\Lambda$ measurements, note first that
  by~\eqref{eq:lambda-consistent} and~\eqref{eq:ld-self-simul} it follows that
  on average over $(x_\alice,x_\bob)\sim \mu_\sampler$,
  \begin{equation}
    \Lambda^{{x}_\alice, {x}_\bob}_{[\eval_{z}(\cdot) = ( \alpha, \beta)]} \ot I
    \simeq_{\delta_3} I \ot \Lambda^{{x}_\alice, {x}_\bob}_{[\eval_{z}(\cdot)
      = ( \alpha, \beta)]} \;.
    \label{eq:lambda-consistent-b}
  \end{equation}
	Using that all outcomes of $\Lambda^{{x}_\alice, {x}_\bob}$ are individual degree $d$
  polynomials and the Schwartz-Zippel lemma (\Cref{lem:schwartz-zippel}) it
  follows that whenever a measurement of $\Lambda^{{x}_\alice, {x}_\bob}\ot
  \Lambda^{{x}_\alice, {x}_\bob}$ returns distinct outcomes, the outcomes take a
  different value at a uniformly random $z$ with probability at least $1-m'd/q$.
  It then follows from~\eqref{eq:lambda-consistent-b} and the fact that
  $\delta_3 \geq m'd/q$ by definition that the strategy $\strategy'$ is accepted
  in Step 1 with probability $1 - O(\delta_3)$.

  Next, the strategy is also accepted in the consistency check performed in item
  2 due to \eqref{eq:ld-strat-con-def}, and the consistency check in item 4(a)
  for the same reasons as for item 1.
  Finally, for the low-degree tests performed in item 3 and items 4(b) and 4(c),
  the strategy succeeds due to consistency and the fact that, as long as both
  players obtain the same polynomial outcomes, they pass the low-degree tests
  with probability 1.

  It remains to analyze the strategy's success probability
  in Step ~\ref{enu:ar-game}, the game check. Recall that this check
  is nontrivial only when $\tvar_{Q,w} = \oracle$ and $\tvar_{\Pi,w} =
  \Point_6$. By assumption the original strategy $\strategy$ succeeds with
  probability $1-O(\eps)$ in the game check. Thus, it will suffice to
  show that the answer distribution from $\strategy'$ is close to the
  answer distribtion from $\strategy$ whenever $\tvar_{Q,w} = \oracle$
  and $\tvar_{\Pi,w} = \Point_6$ for either or both players $w \in
  \{\alice, \bob\}$. For these question types, the measurement
  operator used by $\strategy$ is $M(x_\alice, x_\bob)^{\Point_6,
    z}_{(\alpha, \beta)}$, and
  the operator used by $\strategy'$ is $\Lambda^{(x_\alice,
    x_\bob)}_{[\eval_z(\cdot) = (\alpha, \beta)]}$. By converting
  \Cref{eq:m-lambda-con} and \Cref{eq:lambda-consistent-b} to
  $\approx$ bounds using \Cref{fact:agreement}, and then applying the
  triangle inequality (\Cref{fact:triangle}) we obtain
  that on average over $(x_\alice, x_\bob) \sim \mu_\sampler$,
  \[ \Lambda^{x_\alice, x_\bob}_{[\eval_{z}(\cdot) = (\alpha, \beta)]}
    \ot I \approx_{\delta_3} M(x_\alice,
    x_\bob)^{\Point_6, z}_{(\alpha, \beta)} \ot I_{\bob}. \]
  Thus, restricting to the game check, $\strategy$ and $\strategy'$
  are $\delta_3$-close in the sense of
  \Cref{def:strategy-distance}. They both use the same state, and
  moreover, $\strategy$ is projective. Therefore, by \Cref{lem:close-strategies-have-close-values}, this implies that $\strategy'$ must
  succeed in the game check with probability $1 - \eps -
  O(\delta_3^{1/2}) = 1 - O(\delta_3^{1/2})$.  
  Thus, we have shown that $\strategy'$ succeeds in all subtests of the
  game with probability $1 - O(\delta_3^{1/2})$, and thus the lemma follows.
\end{proof}

We now complete the proof by a reduction to the game $\verifier_n^\ora$: from
the strategy $\strategy'$ we construct a symmetric strategy $\strategy'' =
(\ket{\psi}, A)$ for $\tvora$ by ``decoding'' the low-degree
measurements $G$ and $\Lambda$.
The state $\ket{\psi}$ in $\strategy''$ is identical to the state used in
$\strategy'$ (which is identical to the state used in $\strategy$).
To begin, we define a decoding map $\Delta(\cdot)$, which takes in a polynomial
$g:\F_q^m \to \F_q$ and outputs a string in $\{0,1\}^*$.
This map is computed as follows:
\begin{itemize}
\item First, compute $a = \coded(g) \in \{0,1\}^M$, where $\coded$ is the
  (boolean) decoding of the low-degree code defined in \Cref{sec:ld-encoding}.
\item If there exists an $a_{\mathrm{prefix}} \in \{0, 1\}^*$ of length $\ell_a \leq T$
  such that $a = \mathrm{enc}_\Gamma(a_{\mathrm{prefix}}, \sqcup^{M/2 - \ell_a})$,
  then $\Delta(g) = a_{\mathrm{prefix}}$.
  Otherwise, $\Delta(g)$ is allowed to be arbitrary.
\end{itemize}
We can now define the ``decoded'' measurements $\{A^{{x}_v,v}\}$ and
$\{{A}^{{x}_\alice,{x}_\bob}\}$ as follows:
\begin{equation}
  \label{eq:decoded-meas}
  A^{{x}_v,v}_a  = G^{{x}_v, v}_{[\Delta(\cdot)=a]} \;,\quad
  A^{{x}_\alice,{x}_\bob}_{a_\alice, a_\bob}  =
  \Lambda^{{x}_\alice,{x}_\bob}_{[\Delta(\cdot)_{\alice,\bob}=(a_\alice,a_\bob)]}\;.
\end{equation}
  


\begin{lemma}\label{lem:ar-ora}
  The strategy $\strategy''$ succeeds with probability $1-O(\delta_3^{1/2})$ in the
  game $\tvora_n$.
\end{lemma}
	
\begin{proof}
  We consider the different subtests executed by $\tdora$ (see
  Figure~\ref{fig:oracle-decider}).
  We start with Step~\ref{enu:oracle-full-consistency}, the consistency checks.
  Success in the first check, Step~\ref{enu:check-same}, follows from the success of
  $\strategy'$ in the global consistency check,
  Step~\ref{enu:ar-global-consistency} of $\tdans$, the
  definition~\eqref{eq:decoded-meas}, and the fact that conditioned on
  $\tvar_{Q,\alice}=\tvar_{Q,\bob}=\oracle$, the distribution of
  $(x_{Q,\alice},x_{Q,\bob})$ in $\tvans_n$ is the same as the
  distribution of $(x_{\alice},x_{\bob})$ in $\verifier_n^\ora$, conditioned on
  $\tvar_\alice = \tvar_\bob = \oracle$.
  Similarly, success in the second check, Step~\ref{enu:oracle-versus-player},
  follows from success of $\strategy'$ in the input consistency check,
  Step~\ref{enu:ar-input-consistency} of $\tdans$.

  Next we consider the game check of $\tdora$,
  Step~\ref{enu:oracle-game}.
  To analyze the success probability of $\strategy''$ we use that $\strategy'$
  succeeds in the game check of $\tdans$, Step~\ref{enu:ar-game}, and the
  soundness of the PCP, as shown in Theorem~\ref{thm:pcp-decider}.
  Let $p_{\soundness}$ be as in Theorem~\ref{thm:pcp-decider}.

  Let $w\in\{\alice,\bob\}$, $(x_\alice,x_\bob)$ be in the support of
  $\mu_\sampler$, and $\Pi = (g_1,\ldots,g_5,c_0,\ldots,c_{m'})$ an outcome of
  $\Lambda^{x_\alice,x_\bob}$ such that conditioned on that outcome being
  obtained by player $w$ in the game check of $\tdans$, $\pcpverifier$
  accepts the pair of inputs $(\decider, n , T, \qlen, \gamma, x_{\alice}, x_{\bob})$
  and $(z, a_{w} )$ with probability at least $p_{\soundness}$ over the choice of
  a uniformly random $z\in \F_q^{m'}$ and $a_w = \eval_z(\Pi)$.
			
  For any such $\Pi$, the soundness statement of Theorem~\ref{thm:pcp-decider}
  states that there exist $a_\alice, a_\bob \in \{0,1\}^*$ such that
  $\decider(n, {x}_\alice, {x}_\bob, a_\alice, a_\bob) = 1$ and furthermore
  $a_v = \Delta(g_v)$ for $v \in \{\alice,\bob\}$.   
It follows that for any proof $\Pi$ returned by
  $\{\Lambda^{{x}_\alice,{x}_\bob}_\Pi\}$ which is accepted with probability
  greater than $p_{\soundness}$ in the game check of $\tdans$ it holds
  that $\decider(n, {x}_\alice, {x}_\bob, \Delta(g_\alice), \Delta(g_\bob)) =
  1$.
  Using this observation we evaluate the probability $q_g''$ that the strategy
  $\strategy''$ succeeds in the game check of $\tvora$.
  Let $q_g'$ be the probability that $\strategy'$ succeeds in the game check of
  $\tdans$.
  \begin{align*}
    q_g'' &= \Es{({x}_\alice,{x}_\bob) \sim \mu_{\sampler}} \sum_{a_\alice,
            a_\bob: \decider(n,{x}_\alice,{x}_\bob, a_\alice, a_\bob) = 1}
            \bra{\psi} A^{({x}_\alice,{x}_\bob)}_{a_\alice, a_\bob} \ot I
            \ket{\psi}    \\
          &= \Es{({x}_\alice,{x}_\bob) \sim \mu_{\sampler}} \sum_{\Pi
            : \decider(n,{x}_\alice,{x}_\bob, \Delta(g_1), \Delta(g_2)) = 1}
            \bra{\psi} \Lambda^{({x}_\alice,{x}_\bob)}_{\Pi} \ot I
            \ket{\psi} \\
          &\geq \Es{({x}_\alice,{x}_\bob) \sim \mu_{\sampler}} \sum_{\Pi
            : \decider(n,{x}_\alice,{x}_\bob, \Delta(g_1), \Delta(g_2)) = 1}
            \bra{\psi} \Lambda^{({x}_\alice,{x}_\bob)}_{\Pi} \ot I
            \ket{\psi} \cdot \Pr_{z \sim \F_q^{m'}} [
            \pcpverifier(z,\eval_z(\Pi)) = 1]\\
          &= q_g' -  \Es{({x}_\alice,{x}_\bob) \sim \mu_{\sampler}} \sum_{\Pi
            : \decider(n,{x}_\alice,{x}_\bob, \Delta(g_1), \Delta(g_2)) = 0}
            \bra{\psi} \Lambda^{({x}_\alice,{x}_\bob)}_{\Pi} \ot I
            \ket{\psi} \cdot \Pr_{z \sim \F_q^{m'}} [
            \pcpverifier(z,\eval_z(\Pi))= 1] \\
          &\geq q_g' - ( 1- q_g'')\cdot
            p_{\soundness} \;.
  \end{align*}
  Rearranging terms,
  \begin{equation}
    q_g'' \geq \frac{q_g' - p_{\soundness}}{1 - p_{\soundness}} 
    = 1 - \frac{1-q_g'}{1 - p_{\soundness}}\;. \label{eq:pgameora}
  \end{equation}
  Altogether, using Lemma~\ref{lem:ar-ar} we have shown that $\strategy''$ is
  accepted in each subtest performed by $\tdora$ with probability at
  least $1 - O(\delta_3^{1/2})$.
  Since every subtest occurs with constant probability, the lemma follows.
  \end{proof}
	
	To conclude the proof of the soundness, we appeal to the soundness statement
  for $\tvora$, given in Theorem~\ref{thm:oracle-soundness}.
  Now we establish the form of the error function $\delta$ as stated in the
  theorem. 
  
  Let $a',b'$ be the universal constants from \Cref{lem:ld-soundness} such that
  \[
    \delta_\ld(\eps,q,m,d,t) = a' (dmt)^{a'} (\eps^{b'} + q^{-b'} + 2^{-b' m d}).
 	\]
	By our choice of PCP parameters in \Cref{def:pcpparams}, we have that $m \leq
  s$, so therefore we can upper bound the product $dm'(m' + 6)$ by $c \gamma
  s^3$ where $c > 1$ is a universal constant.
  Furthermore, we see from \Cref{def:pcpparams} that $2^{dm} \geq q \geq s^{(\gamma
    b' + 3a')/b'}$, so we have
	\begin{align}
		\delta_\ld(\eps,q,m',d,m'+6)
    & \leq a' (c\gamma)^{a'} ( s^{3a'} \cdot \eps^{b'} +
      2 \cdot s^{3 a'} \cdot q^{-b'}) \notag \\
    & \leq 2 a' (c\gamma)^{a'} ( s^{3a'} \cdot \eps^{b'} + s^{3 a'}
      \cdot s^{-(\gamma b' + 3a')}) \notag \\
    & \leq a'' \gamma^{a''} ( s^{a''} \cdot \eps^{b''}
      + s^{-\gamma b''}) \label{eq:ar-sound-3}
	\end{align}
	where we set $a'' = \max \{ 2 a' c^{a'} , 3a'\}$ and $b'' = b'$.
From \Cref{prop:explicit-padded-succinct-deciders,eq:ar-time-assumption} we
  have that $s$ is at most $\poly( (\lambda n)^\mu , \sigma)$, so there exists a
  universal constant $C > 0$ such that for all $n \geq 2$,
  \[
  	s \leq C ( \lambda \cdot \sigma \cdot n)^{C \mu}.
  \]
  Plugging this bound into \Cref{eq:ar-sound-3} and using the choice of $q,d,m'$
  from \Cref{def:pcpparams}, we get that there exist universal constants $a > 1$
  and $0 < b < 1$ such that
  \begin{align*}
    \delta_3^{1/2}
    & = O \left ( \Big (\delta_\ld(\eps, q, m', d, m'+6)^{1/2} +
      \frac{m'd}{q} \Big)^{1/4} \right ) \\
    & \leq a \gamma^a ( (\lambda \cdot \sigma \cdot n)^{\mu a} \cdot \eps^{b}
      + (\lambda \cdot \sigma \cdot n)^{- \mu b \gamma}).
  \end{align*}  
  


  To show the entanglement bound, we observe that the strategy $\strategy''$ for
  $\tvora$ constructed above uses the same entangled state $\ket{\psi}$ as the
  strategy $\strategy$ for $\tvans$ that we started with; the claimed bound
  follows.
\end{proof}


 
 
\section{Parallel Repetition}
